\section{Hệ bất phương trình bậc nhất hai ẩn}

\subsection{Đề 1}
\Opensolutionfile{ans}[ans/ans-0-B1]
\caulc

\begin{ex}%[0D2H2-2]
	Trong các điểm sau đây, điểm nào thuộc miền nghiệm của hệ bất phương trình $\heva{&x+3y-2 \ge 0\\&2x+y+1 \le 0}$.		
	\choice{$\left(0; 1\right)$}
	{\True $\left(-1; 1\right)$}
	{$\left(1; 3\right)$}
	{$\left(-1; 0\right)$}
	\loigiai{
		 \begin{itemize}
		 	\item Thay $(0; 1)$ vào hệ ta được $\heva{&0+3-2>0\\&0+1+1>0}$. Do đó điểm $(0;1)$ không thoả mãn.
		 	\item Thay $(-1; 1)$ vào hệ ta được $\heva{&-1+3-2=0\\&-2+1+1=0}$. Do đó điểm $(0;1)$  thoả mãn.
		 	\item Thay $(1; 3)$ vào hệ ta được $\heva{&1+9-2>0\\&2+1+1>0}$. Do đó điểm $(1;3)$ không thoả mãn.
		 	\item Thay $(-1; 0)$ vào hệ ta được $\heva{&-1+0-2<0\\&-2+1+1=0}$. Do đó điểm $(0;1)$ không thoả mãn.
		 \end{itemize}
	}
\end{ex}

\begin{ex}%[0D2H2-2]
	Cho miền gạch chéo như hình vẽ dưới đây
	\begin{center}
		\begin{tikzpicture}[scale=0.7, line join=round, line cap=round, >=stealth]
			\tikzset{every node/.style={scale=1}}
			\def\xmin{-3}\def\xmax{2.5}\def\ymin{-2}\def\ymax{4}
			\draw[->] (\xmin-0.2,0)--(\xmax+0.2,0) node[below]{$x$};
			\draw[->] (0,\ymin-0.2)--(0,\ymax+0.2) node[right]{$y$};
			\draw (0,0) node[above left]{$O$};
			\foreach \x in {-1}\draw (\x,0.1)--(\x,-0.1) node[below]{$\x$};
			\foreach \y in {3}\draw (0.1,\y)--(-0.1,\y) node[right]{$\y$};
			\clip (\xmin,\ymin) rectangle (\xmax,\ymax);
			
			\draw[thick,smooth,samples=200,domain=\xmin:\xmax] plot (\x,{-2*(\x)+1});
			\clip (\xmin,\ymin) rectangle (\xmax,\ymax);
			\path
			(-1.5,3.3) node[above left]{$d_1$}
			(2,2) node[above left]{$d_2$}
			(-2,-1) node[below]{$d_3$};
			\draw[thick,smooth,samples=200,domain=\xmin:\xmax] plot (\x,{0.5*(\x)+1});
			
			\draw[thick,smooth,samples=200,domain=\xmin:\xmax] plot (\x,{3*(\x)+6});
			\draw[dashed] (0,3) -- (-1,3) -- (-1,0);
			
			\fill[pattern=north east lines,smooth]
			plot[domain=-2:-1] (\x,{3*(\x)+6})
			-- plot[domain=-1:-2] (\x,{0.5*(\x)+1})
			-- cycle;
			
			\fill[pattern=north east lines,smooth]
			plot[domain=-1:0] (\x,{-2*(\x)+1})
			-- plot[domain=0:-1] (\x,{0.5*(\x)+1})
			-- cycle;			
		\end{tikzpicture}
	\end{center}
	Miền trên đây biểu diễn tập nghiệm của hệ bất phương trình nào?
	\choice{$\heva{&2x+y>1\\&-x+2y<2\\&3x-y>6}$}
	{ $\heva{&2x+y<1\\&-x+2y<2\\&3x-y>-6}$}
	{\True$\heva{&2x+y<1\\&-x+2y>2\\&3x-y>-6}$}
	{$\heva{&2x+y>1\\&x-2y<2\\&3x-y>6}$}
	\loigiai{Lấy điểm $B\left(-1 ; 1\right)$ thuộc miền gạch chéo.
		\begin{itemize}
			\item Thay $(-1,1)$ vào hệ $\heva{&2x+y>1\\&-x+2y<2\\&3x-y>6}$ ta được $\heva{&-2+1<1\\&1+2>2\\&-3-1<6}$ không thoả mãn. Do đó $\heva{&2x+y>1\\&-x+2y<2\\&3x-y>6}$ không thoả mãn.
			\item Thay $(-1,1)$ vào hệ $\heva{&2x+y<1\\&-x+2y<2\\&3x-y>-6}$ ta được $\heva{&-2+1<1\\&1+2>2\\&-3-1<-6}$. Do đó $\heva{&2x+y<1\\&-x+2y<2\\&3x-y>-6}$ không thoả mãn.
			\item Thay $(-1,1)$ vào hệ $\heva{&2x+y<1\\&-x+2y<2\\&3x-y>-6}$ ta được $\heva{&-2+1<1\\&1+2>2\\&-3-1>-6}$. Do đó $\heva{&2x+y<1\\&-x+2y<2\\&3x-y>-6}$ không thoả mãn.
			\item Thay $(-1,1)$ vào hệ $\heva{&2x+y<1\\&-x+2y>2\\&3x-y>-6}$ ta được $\heva{&-2+1<1\\&1+2>2\\&-3-1>-6}$. Do đó $\heva{&2x+y<1\\&-x+2y<2\\&3x-y>-6}$ thoả mãn.
			\item Thay $(-1,1)$ vào hệ $\heva{&2x+y>1\\&-x+2y<2\\&3x-y>6}$ ta được $\heva{&-2+1<1\\&1+2>2\\&-3-1<6}$. Do đó $\heva{&2x+y>1\\&-x+2y<2\\&3x-y>6}$ không thoả mãn.
		\end{itemize} }
	\end{ex}
	
	\begin{ex}%[0D2H2-2]
		Miền nghiệm của hệ bất phương trình $\heva{&x-y>0\\&x-3y \le -3\\&x+y>5}$ \textbf{không chứa} điểm nào sau đây?
		\choice{\True $A\left(3; 2\right)$}
		{$B\left(6; 3\right)$}
		{$C\left(6; 4\right)$}
		{$D\left(5; 4\right)$}
		\loigiai{
			Xét điểm $A \left(3; 2\right)$ ta có $\heva{&3-2>0\\&3-3\cdot 2 \le -3\\&3+2>5}$ không thỏa mãn. Do đó mền nghiệm của hệ bất phương trình $\heva{&x-y>0\\&x-3y \le -3\\&x+y>5}$ \textbf{không chứa} điểm $A\left(3; 2\right)$.
		}
	\end{ex}
	
	\begin{ex}%[0D2H2-2]
		Miền nghiệm của hệ bất phương trình $\heva{&3-y < 0\\&2x-3y+1 > 0}$ chứa điểm nào sau đây?
		\choice{$\left(3; 4\right)$}
		{$\left(4; 3\right)$}
		{\True $\left(7; 4\right)$}
		{$\left(4; 4\right)$}
		\loigiai{
		Trước hết, ta vẽ hai đường thẳng:\\
			$ \left(d_1\right): 3-y=0 \\
			\left(d_2\right): 2 x-3 y+1=0$
			\begin{center}
			\begin{tikzpicture}[scale=0.7, line join=round, line cap=round, >=stealth]
				\tikzset{every node/.style={scale=1}}
				\def\xmin{-1}\def\xmax{9}\def\ymin{-1}\def\ymax{6}
				\draw[->] (\xmin-0.2,0)--(\xmax+0.2,0) node[below]{$x$};
				\draw[->] (0,\ymin-0.2)--(0,\ymax+0.2) node[right]{$y$};
				\draw (0,0) node[below right]{$O$};
				\foreach \x in {4}\draw (\x,0.1)--(\x,-0.1) node[below]{$\x$};
				\foreach \y in {3}\draw (0.1,\y)--(-0.1,\y) node[above left]{$\y$};
				\clip (\xmin,\ymin) rectangle (\xmax,\ymax);
				
				\draw[thick,smooth,samples=200,domain=\xmin:\xmax] plot (\x,{0*(\x)+3});
				\clip (\xmin,\ymin) rectangle (\xmax,\ymax);
				
				\draw[thick,smooth,samples=200,domain=\xmin:\xmax] plot (\x,{(2/3)*(\x)+(1/3)});
				
				\fill[pattern=north east lines,smooth]
				plot[domain=-1:9] (\x,{0*(\x)+3})
				-- plot[domain=9:-1] (\x,{0*(\x)-1})
				-- cycle;
				
				\fill[pattern=north east lines,smooth]
				plot[domain=-1:9] (\x,{0*(\x)+6})
				-- plot[domain=9:-1] (\x,{(2/3)*(\x)+(1/3)})
				-- cycle;
				
				\path
				(8,4) node[above left]{$d_1$}
				(8,3) node[above left]{$d_1$};
				\end{tikzpicture}
			\end{center}
		Ta thấy $\left(6 ; 4\right)$ là nghiệm của hai bất phương trình. Điều đó có nghĩa điểm $\left(6 ; 4\right)$ thuộc cả hai miền nghiệm của hai bất phương trình. Sau khi gạch bỏ các miền không thích hợp, miền không bị gạch là miền nghiệm của hệ.
		}
	\end{ex}
	
	\begin{ex}%[0D2H2-2]
		Giá trị nhỏ nhất của biết thức $F=y-x$ trên miền xác định bởi hệ: $\heva{&y-2x \le 2 \\ & 2y-x \ge 4 \\ & x+5 \le 5}$ là:
		\choice{\True $\min F=1$ khi $x=2, y=3$}
		{$\min F=2 k h i x=0, y=2$}
		{$\min F=3$ khi $x=1, y=4$}
		{Không tồn tại giá trị nhỏ nhất của $F$}
		\loigiai{
			Ta tìm giao điểm của các cặp đường thẳng trong miền xác định của hệ $\heva{& y-2x \le 2 \\ &2y-x \ge 4 \\ &x+y \le 5}$
			
			$\heva{& -2x+y=2 \\ & -x+2y=4} \Leftrightarrow \heva{& x=0 \\ &y=2} \Rightarrow A\left(0;2\right)$
			
			$\heva{& -2x+y=2 \\ & x+y=5} \Leftrightarrow \heva{& x=\dfrac{3}{2} \\ &y=\dfrac{7}{2}} \Rightarrow B\left(\dfrac{3}{2};\dfrac{7}{2}\right)$\\
			
			$\heva{& x+y=5 \\ & -x+2y=4} \Leftrightarrow \heva{& x=2  &y=3} \Rightarrow C\left(2;3\right)$
			
			Ta tính giá trị của $F=y-x$ tại các giao điểm:\\
			
			Tại $A\left(0 ; 2\right) \Rightarrow F=2-0=2$.\\
			
			Tại $B\left(\dfrac{3}{2} ; \dfrac{7}{2}\right) \Rightarrow F=\dfrac{7}{2}-\dfrac{3}{2}=2$.
			
			Tại $C\left(2 ; 3\right) \Rightarrow F=3-2=1$.
			
			Vậy $\min F=1$ khi $x=2, y=3$}
	\end{ex}
	
	\begin{ex}%[0D2H2-2]
		Biểu thức $F=y-x$ đạt giá trị nhỏ nhất với điều kiện $\heva{&-2x+y \le -2 \\ &x-2y \le 2 \\ &x+y \le 5 \\ &x \ge 0}$ tại điểm $S\left(x; y\right)$ có toạ độ là
		\choice{\True $\left(4; 1\right)$}
		{$\left(3; 1\right)$}
		{$\left(2; 1\right)$}
		{$\left(1; 1\right)$}
		\loigiai{
			\textbf{Cách 1: Thử máy tính.} Ta dùng máy tính lần lượt kiểm tra các đáp án để xem đáp án nào thỏa hệ bất phương trình trên loại được đáp án $D$.
			
			Ta lần lượt tính hiệu $F=y-x$ và $\min F=-3$ tại $x=4, y=1$.
			
			\textbf{Cách 2: Tự luận:}
	\begin{center}
		\begin{tikzpicture}[scale=0.7, line join=round, line cap=round, >=stealth]
			\tikzset{every node/.style={scale=1}}
			\def\xmin{-3}\def\xmax{8}\def\ymin{-3}\def\ymax{6}
			\draw[->] (\xmin-0.2,0)--(\xmax+0.2,0) node[below]{$x$};
			\draw[->] (0,\ymin-0.2)--(0,\ymax+0.2) node[right]{$y$};
			\draw (0,0) node[below left]{$O$};
			\foreach \x in {}\draw (\x,0.1)--(\x,-0.1) node[below]{$\x$};
			\foreach \y in {}\draw (0.1,\y)--(-0.1,\y) node[above left]{$\y$};
			\clip (\xmin,\ymin) rectangle (\xmax,\ymax);
			
			\draw[thick,smooth,samples=200,domain=\xmin:\xmax] plot (\x,{2*(\x)-2});
			\clip (\xmin,\ymin) rectangle (\xmax,\ymax);
			
			\draw[thick,smooth,samples=200,domain=\xmin:\xmax] plot (\x,{(1/2)*(\x)-1});
			
			\draw[thick,smooth,samples=200,domain=\xmin:\xmax] plot (\x,{-1*(\x)+5});
			
			\fill[pattern=north east lines,smooth]
			plot[domain=-3:8] (\x,{0*(\x)+6})
			-- plot[domain=8:-3] (\x,{2*(\x)-2})
			-- cycle;
			
			\fill[pattern=north east lines,smooth]
			plot[domain=-3:8] (\x,{(1/2)*(\x)-1})
			-- plot[domain=8:-3] (\x,{0*(\x)-3})
			-- cycle;
			
			\fill[pattern=north east lines,smooth]
			plot[domain=-3:8] (\x,{0*(\x)+6})
			-- plot[domain=8:-3] (\x,{-1*(\x)+5})
			-- cycle;
			
			\path
			(7/3,8/3) node[below]{$A$}
			(2/3,-2/3) node[below right]{$B$}
			(4,1) node[above]{$C$};
			
		\end{tikzpicture}
	\end{center}
			Tọa độ $A\left(\dfrac{7}{3}; \dfrac{8}{3}\right), B\left(\dfrac{2}{3};-\dfrac{2}{3}\right), C\left(4; 1\right)$. Giá trị $F$ lần lượt tại toạ độ các điểm $B, C, A$ là $-\dfrac{4}{3},-3; \dfrac{1}{3}$. 
			
			Suy ra $\min F=-3$ tại $\left(4; 1\right)$.
		}
	\end{ex}
	
	\begin{ex}%[0D2H2-2]
		Giá trị nhỏ nhất của biều thức $F=y-x$ trên miền xác định bởi hệ $\heva{&y-2x \le 2 \\ &2y-x \ge 4 \\ &x+5 \le 5}$ là
		\choice{\True $\min F=1$ khi $x=2, y=3$}
		{$\min F=2$ khi $x=0, y=2$}
		{$\min F=3$ khi $x=1, y=4$}
		{$\min F=0$ khi $x=0, y=0$}
		\loigiai{
			Miền nghiệm của hệ $\heva{&y-2x \le 2 \\ &2y-x \ge 4 \\ &x+5 \le 5}$ là miền trong của tam giác $ABC$ kể cả biên (như hình).
	\begin{center}
		\begin{tikzpicture}[scale=0.7, line join=round, line cap=round, >=stealth]
			\tikzset{every node/.style={scale=1}}
			\def\xmin{-3}\def\xmax{5}\def\ymin{-2}\def\ymax{6}
			\draw[->] (\xmin-0.2,0)--(\xmax+0.2,0) node[below]{$x$};
			\draw[->] (0,\ymin-0.2)--(0,\ymax+0.2) node[right]{$y$};
			\draw (0,0) node[below left]{$O$};
			\foreach \x in {}\draw (\x,0.1)--(\x,-0.1) node[below]{$\x$};
			\foreach \y in {}\draw (0.1,\y)--(-0.1,\y) node[above left]{$\y$};
			\clip (\xmin,\ymin) rectangle (\xmax,\ymax);
			
			\draw[thick,smooth,samples=200,domain=\xmin:\xmax] plot (\x,{2*(\x)+2});
			\clip (\xmin,\ymin) rectangle (\xmax,\ymax);
			
			\draw[thick,smooth,samples=200,domain=\xmin:\xmax] plot (\x,{(1/2)*(\x)+2});
			
			\draw[thick,smooth,samples=200,domain=\xmin:\xmax] plot (\x,{-1*(\x)+5});
			
			\fill[pattern=north east lines,smooth]
			plot[domain=-3:6] (\x,{0*(\x)+6})
			-- plot[domain=6:-3] (\x,{2*(\x)+2})
			-- cycle;
			
			\fill[pattern=north east lines,smooth]
			plot[domain=-2:6] (\x,{(1/2)*(\x)+2})
			-- plot[domain=6:-2] (\x,{0*(\x)-3})
			-- cycle;
			
			\fill[pattern=north east lines,smooth]
			plot[domain=-3:8] (\x,{0*(\x)+6})
			-- plot[domain=8:-3] (\x,{-1*(\x)+5})
			-- cycle;
			
			\path
			(0,2) node[below right]{$A$}
			(1,4) node[right]{$B$}
			(2,3) node[below]{$C$};
			
		\end{tikzpicture}
	\end{center}
			Ta thấy $F=y-x$ đạt giá trị nhỏ nhất chỉ có thể tại các điểm $A, B, C$.
			
			Tại $A\left(0; 2\right)$ thì $F=2$.
			
			Tại $B\left(1; 4\right)$ thì $F=3$
			
			Tại $A\left(2; 3\right)$ thì $F=1$.
			
			Vậy $\min F=1$ khi $x=2, y=3$.
		}
	\end{ex}
	
	\begin{ex}%[0D2H2-2]
		Giá trị nhỏ nhất của biết thức $F=y-x$ trên miền xác định bởi hệ $\heva{&2x+y \le 2 \\ &x-y \le 2 \\ &5x+y \ge -4}$ là
		\choice{$\min F=-3$ khi $x=1, y=-2$}
		{$\min F=0$ khi $x=0, y=0$}
		{\True $\min F=-2$ khi $x=\dfrac{4}{3}, y=-\dfrac{2}{3}$}
		{$\min F=8$ khi $x=-2, y=6$}
		\loigiai{
		Biểu diễn miền nghiệm của hệ bất phương trình $\heva{&2x+y \le 2 \\ &x-y \le 2 \\ &5x+y \ge -4}$ trên hệ trục tọa độ như dưới đây:
	\begin{center}
		\begin{tikzpicture}[scale=0.7, line join=round, line cap=round, >=stealth]
			\tikzset{every node/.style={scale=1}}
			\def\xmin{-4}\def\xmax{5}\def\ymin{-4}\def\ymax{7}
			\draw[->] (\xmin-0.2,0)--(\xmax+0.2,0) node[below]{$x$};
			\draw[->] (0,\ymin-0.2)--(0,\ymax+0.2) node[right]{$y$};
			\draw (0,0) node[below left]{$O$};
			\foreach \x in {}\draw (\x,0.1)--(\x,-0.1) node[below]{$\x$};
			\foreach \y in {}\draw (0.1,\y)--(-0.1,\y) node[above left]{$\y$};
			\clip (\xmin,\ymin) rectangle (\xmax,\ymax);
			
			\draw[thick,smooth,samples=200,domain=\xmin:\xmax] plot (\x,{-2*(\x)+2});
			\clip (\xmin,\ymin) rectangle (\xmax,\ymax);
			
			\draw[thick,smooth,samples=200,domain=\xmin:\xmax] plot (\x,{1*(\x)-2});
			
			\draw[thick,smooth,samples=200,domain=\xmin:\xmax] plot (\x,{-5*(\x)-4});
			
			\fill[pattern=north east lines,smooth]
			plot[domain=-4:5] (\x,{0*(\x)+7})
			-- plot[domain=5:-4] (\x,{-2*(\x)+2})
			-- cycle;
			
			\fill[pattern=north east lines,smooth]
			plot[domain=-4:5] (\x,{1*(\x)-2})
			-- plot[domain=5:-4] (\x,{0*(\x)-4})
			-- cycle;
			
			\fill[pattern=north east lines,smooth]
			plot[domain=-4:5] (\x,{-5*(\x)-4})
			-- plot[domain=5:-4] (\x,{0*(\x)-4})
			-- cycle;
			
			\path
			(-2,6) node[below right]{$A$}
			(4/3,-2/3) node[right]{$B$}
			(-1/3,-7/3) node[left]{$C$};
			
		\end{tikzpicture}
	\end{center}
		Giá trị nhỏ nhất của biết thức $F=y-x$ chỉ đạt được tại các điểm
		
		$A\left(-2;6\right), C\left(\dfrac{4}{3};-\dfrac{2}{3}\right), B\left(-\dfrac{1}{3};-\dfrac{7}{3}\right)$
			
		Ta có: $F\left(A\right)=8; F\left(B\right)=-2; F\left(C\right)=-2$.
			
		Vậy $\min F=-2$ khi $x=\dfrac{4}{3}, y=-\dfrac{2}{3}$}
	\end{ex}
	
	\begin{ex}%[0D2H2-2]
		Cho hệ bất phương trình $\heva{&x-y \le 2 \\ &3x+5y \le 15 \\ &x \ge 0 \\ &y \ge 0}$. Khẳng định nào sau đây là khẳng định sai?
		\choice{Trên mặt phẳng tọa độ $Ox y$, biểu diễn miền nghiệm của hệ bất phương trình đã cho là miền tứ giác $ABCO$ kể cả các cạnh với $A\left(0; 3\right), B\left(\dfrac{25}{8}; \dfrac{9}{8}\right), C\left(2; 0\right)$ và $O\left(0; 0\right)$}
		{\True Đường thẳng $\Delta: x+y=m$ có giao điểm với tứ giác $ABCO$ kể cả khi $-1 \leq m \leq \dfrac{17}{4}$}
		{Giá trị lớn nhất của biểu thức $x+y$, với $x$ và $y$ thỏa mãn hệ bất phương trình đã cho là $\dfrac{17}{4}$}
		{Giá trị nhỏ nhất của biểu thức $x+y$, với $x$ và $y$ thỏa mãn hệ bất phương trình đã cho là 0}
		\loigiai{Trước hết, ta vẽ bốn đường thẳng:
		
		$\left(d_{1}: x-y=2\right)$
		
		$\left(d_{2}: 3x+5y=15\right)$
		
		$\left(d_{3}: x=0\right)$
		
		$\left(d_{4}: y=0\right)$
	\begin{center}
		\begin{tikzpicture}[scale=0.7, line join=round, line cap=round, >=stealth]
			\tikzset{every node/.style={scale=1}}
			\def\xmin{-3}\def\xmax{5}\def\ymin{-3}\def\ymax{5}
			\draw[->] (\xmin-0.2,0)--(\xmax+0.2,0) node[below]{$x$};
			\draw[->] (0,\ymin-0.2)--(0,\ymax+0.2) node[right]{$y$};
			\draw (0,0) node[below left]{$O$};
			\foreach \x in {}\draw (\x,0.1)--(\x,-0.1) node[below]{$\x$};
			\foreach \y in {}\draw (0.1,\y)--(-0.1,\y) node[above left]{$\y$};
			\clip (\xmin,\ymin) rectangle (\xmax,\ymax);
			
			\draw[thick,smooth,samples=200,domain=\xmin:\xmax] plot (\x,{1*(\x)-2});
			\clip (\xmin,\ymin) rectangle (\xmax,\ymax);
			
			\draw[thick,smooth,samples=200,domain=\xmin:\xmax] plot (\x,{(-3/5)*(\x)+3});
			
			%\draw[thick,smooth,samples=200,domain=\xmin:\xmax] plot (\x,{-5*(\x)-4});
			
			\fill[pattern=north east lines,smooth]
			plot[domain=-4:5] (\x,{1*(\x)-2})
			-- plot[domain=5:-4] (\x,{0*(\x)-3})
			-- cycle;
			
			\fill[pattern=north east lines,smooth]
			plot[domain=-4:5] (\x,{0*(\x)+7})
			-- plot[domain=5:-4] (\x,{(-3/5)*(\x)+3})
			-- cycle;
			
			\fill[pattern=north east lines,smooth]
			plot[domain=-4:0] (\x,{0*(\x)+7})
			-- plot[domain=0:-4] (\x,{0*(\x)-4})
			-- cycle;
			
			\fill[pattern=north east lines,smooth]
			plot[domain=0:5] (\x,{0*(\x)+0})
			-- plot[domain=5:0] (\x,{0*(\x)-4})
			-- cycle;
			
		\end{tikzpicture}
	\end{center}
		Miền nghiệm là phần không bị gạch, kể cả biên.}
	\end{ex}
	
	\begin{ex}%[0D2H2-2]
		Giá trị lớn nhất của biết thức $F(x; y)=x+2 y$ với điều kiện $\heva{&0 \le y \le 4 \\ &x \ge 0 \\ &x-y-1 \le 0 \\ &x+2y-10 \le 0}$ là
		\choice{$6$}
		{$8$}
		{\True $10$}
		{$12$}
		\loigiai{
			Vẽ đường thẳng $d_{1}: x-y-1=0$, đường thẳng $d_{1}$ qua hai điểm $\left(0;-1\right)$ và $\left(1; 0\right)$.\\			
			Vẽ đường thẳng $d_{2}: x+2 y-10=0$, đường thẳng $d_{2}$ qua hai điểm $\left(0; 5\right)$ và $\left(2; 4\right)$.\\			
			Vẽ đường thẳng $d_{3}: y=4$.
	\begin{center}
		\begin{tikzpicture}[scale=0.7, line join=round, line cap=round, >=stealth]
			\tikzset{every node/.style={scale=1}}
			\def\xmin{-2}\def\xmax{6}\def\ymin{-2}\def\ymax{7}
			\draw[->] (\xmin-0.2,0)--(\xmax+0.2,0) node[below]{$x$};
			\draw[->] (0,\ymin-0.2)--(0,\ymax+0.2) node[right]{$y$};
			\draw (0,0) node[below left]{$O$};
			\foreach \x in {}\draw (\x,0.1)--(\x,-0.1) node[below]{$\x$};
			\foreach \y in {}\draw (0.1,\y)--(-0.1,\y) node[above left]{$\y$};
			\clip (\xmin,\ymin) rectangle (\xmax,\ymax);
			
			\draw[thick,smooth,samples=200,domain=\xmin:\xmax] plot (\x,{1*(\x)-1});
			\clip (\xmin,\ymin) rectangle (\xmax,\ymax);
			
			\draw[thick,smooth,samples=200,domain=\xmin:\xmax] plot (\x,{(-1/2)*(\x)+5});
			
			\draw[thick,smooth,samples=200,domain=\xmin:\xmax] plot (\x,{0*(\x)+4});
			
			\fill[pattern=north east lines,smooth]
			plot[domain=-4:7] (\x,{1*(\x)-1})
			-- plot[domain=7:-4] (\x,{0*(\x)-3})
			-- cycle;
			
			\fill[pattern=north east lines,smooth]
			plot[domain=-4:7] (\x,{0*(\x)+7})
			-- plot[domain=7:-4] (\x,{(-1/2)*(\x)+5})
			-- cycle;
			
			\fill[pattern=north east lines,smooth]
			plot[domain=-2:7] (\x,{0*(\x)+7})
			-- plot[domain=7:-2] (\x,{0*(\x)+4})
			-- cycle;
			
			\fill[pattern=north east lines,smooth]
			plot[domain=-2:7] (\x,{0*(\x)+0})
			-- plot[domain=7:-2] (\x,{0*(\x)-4})
			-- cycle;
			
			\fill[pattern=north east lines,smooth]
			plot[domain=-2:0] (\x,{0*(\x)+6})
			-- plot[domain=0:-2] (\x,{0*(\x)-2})
			-- cycle;
			
			\path
			(4,3) node[right]{$A$}
			(2,4) node[below]{$B$}
			(0,4) node[below right]{$C$}
			(1,0) node[above]{$F$}
			;
		\end{tikzpicture}
	\end{center}
			Miền nghiệm là ngũ giác $ABCOE$ với $A\left(4; 3\right), B\left(2; 4\right), C\left(0; 4\right), E\left(1; 0\right)$.\\			
			Ta có: $F\left(4; 3\right)=10, F\left(2; 4\right)=10, F\left(0; 4\right)=8, F\left(1; 0\right)=1, F\left(0; 0\right)=0$.\\			
			Vậy giá trị lớn nhất của biết thức $F\ (x; y)=x+2 y$ bằng 
			$10$.}
	\end{ex}
	
	\begin{ex}%[0D2H2-2]
		Biểu thức $L=y-x$, với $x$ và $y$ thỏa mãn hệ bất phương trình $\heva{&2x+3y-6 \le 0 \\ &x \ge 0 \\ &2x-3y-1 \le 0}$, đạt giá trị lớn nhất là $a$ và đạt giá trị nhỏ nhất là $b$. Hãy chọn kết quả đúng trong các kết quả sau
		\choice{$a=\dfrac{25}{8}$ và $b=-2$}
		{\True $a=2$ và $b=-\dfrac{11}{12}$}
		{$a=3$ và $b=0$}
		{$a=3$ và $b=\dfrac{-9}{8}$}
		\loigiai{
		Trước hết, ta vẽ ba đường thẳng:\\		
		$\left(d_{1}\right): 2 x+3 y-6=0$.\\		
		$\left(d_{2}\right): x=0$.\\		
		$\left(d_{3}\right): 2 x-3 y-1=0$
	\begin{center}
		\begin{tikzpicture}[scale=1, line join=round, line cap=round, >=stealth]
			\tikzset{every node/.style={scale=1}}
			\def\xmin{-2}\def\xmax{3}\def\ymin{-2}\def\ymax{4}
			\draw[->] (\xmin-0.2,0)--(\xmax+0.2,0) node[below]{$x$};
			\draw[->] (0,\ymin-0.2)--(0,\ymax+0.2) node[right]{$y$};
			\draw (0,0) node[below left]{$O$};
			\foreach \x in {}\draw (\x,0.1)--(\x,-0.1) node[below]{$\x$};
			\foreach \y in {}\draw (0.1,\y)--(-0.1,\y) node[above left]{$\y$};
			\clip (\xmin,\ymin) rectangle (\xmax,\ymax);
			
			\draw[thick,smooth,samples=200,domain=\xmin:\xmax] plot (\x,{(-2/3)*(\x)+2});
			\clip (\xmin,\ymin) rectangle (\xmax,\ymax);
			
			\draw[thick,smooth,samples=200,domain=\xmin:\xmax] plot (\x,{(2/3)*(\x)-(1/3)});
			
			%\draw[thick,smooth,samples=200,domain=\xmin:\xmax] plot (\x,{0*(\x)+4});
			
			\fill[pattern=north east lines,smooth]
			plot[domain=-4:7] (\x,{0*(\x)+5})
			-- plot[domain=7:-4] (\x,{(-2/3)*(\x)+2})
			-- cycle;
			
			\fill[pattern=north east lines,smooth]
			plot[domain=-4:7] (\x,{(2/3)*(\x)-(1/3)})
			-- plot[domain=7:-4] (\x,{0*(\x)-5})
			-- cycle;
			
			\fill[pattern=north east lines,smooth]
			plot[domain=-2:7] (\x,{0*(\x)+7})
			-- plot[domain=7:-2] (\x,{0*(\x)+4})
			-- cycle;
			
			\fill[pattern=north east lines,smooth]
			plot[domain=-2:7] (\x,{0*(\x)+0})
			-- plot[domain=7:-2] (\x,{0*(\x)-4})
			-- cycle;
			
			\fill[pattern=north east lines,smooth]
			plot[domain=-2:0] (\x,{0*(\x)+6})
			-- plot[domain=0:-2] (\x,{0*(\x)-2})
			-- cycle;
			
			\path
			(0,2) node[right]{$A$}
			(7/4,5/6) node[below]{$B$}
			(0,-1/3) node[below right]{$C$}
			;
		\end{tikzpicture}
	\end{center}
		Ta thấy $\left(0; 0\right)$ là nghiệm của cả ba bất phương trình. Điều đó có nghĩa gốc tọa độ thuộc cả ba miền nghiệm của cả ba bất phương trình. Sau khi gạch bỏ các miền không thích hợp, miền không bị gạch là miền nghiệm của hệ (kể cả biên).\\			
		Miền nghiệm là hình tam giác $ABC$ (kể cả biên), với $A\left(0; 2\right), B\left(\dfrac{7}{4}; \dfrac{5}{6}\right), C\left(0;-\dfrac{1}{3}\right)$.\\			
		Vậy ta có $a=2-0=2, b=\dfrac{5}{6}-\dfrac{7}{4}=-\dfrac{11}{12}$.
		}
	\end{ex}

	\begin{ex}%[0D2V2-3]
		Trong một cuộc thi pha chế, hai đội $ A$, $B $ được sử dụng tối đa $24$ g hương liệu, $9$ lít nước và $210$ g đường để pha chế nước cam và nước táo. Để pha chế $1$ lít nước cam cần $30$ g đường, $1$ lít nước và $1$ g hương liệu; pha chế $1$ lít nước táo cần $10$ g đường, $1$ lít nước và $4$ g hương liệu. Mỗi lít nước cam nhận được $60$ điểm thưởng, mỗi lít nước táo nhận được $80$ điểm thưởng. Đội $ A $ pha chế được $a$ lít nước cam và $b$ lít nước táo và dành được điểm thưởng cao nhất. Hiệu số $a-b$ là
		\choice{$1$}
		{$3$}
		{\True $-1$}
		{$-6$}
		\loigiai{
			Gọi $x$, $y$ lần lượt là số lít nước cam và nước táo mà mỗi đội cần pha chế $\left(x \geq 0; y \geq 0\right)$.\\			
			Để pha chế $x$ lít nước cam cần $30 x $ g đường, $x$ lít nước và $x$ g hương liệu.\\			
			Để pha chế $y$ lít nước táo cần $10 y$ g đường, $y$ lít nước và $4 y$ g hương liệu.\\			
			Theo bài ra ta có hệ bất phương trình
			$$\heva{&30x+10y \le 210 \\ &x+y \le 9  (*) \\ &x+4y \le 24 \\ &x \ge 0;y \ge 0  }$$ 
			Số điểm đạt được khi pha $x$ lít nước cam và $y$ lít nước táo là $M\left(x, y\right)=60 x+80 y$. Bài toán trở thành tìm $x$, $y$ để $M\left(x, y\right)$ đạt giá trị lớn nhất.\\			
			Ta biểu diễn miền nghiệm của hệ (*) trên mặt phẳng tọa độ như sau
	\begin{center}
		\begin{tikzpicture}[scale=0.7, line join=round, line cap=round, >=stealth]
			\tikzset{every node/.style={scale=1}}
			\def\xmin{-2}\def\xmax{10}\def\ymin{-2}\def\ymax{10}
			\draw[->] (\xmin-0.2,0)--(\xmax+0.2,0) node[below]{$x$};
			\draw[->] (0,\ymin-0.2)--(0,\ymax+0.2) node[right]{$y$};
			\draw (0,0) node[below left]{$O$};
			\foreach \x in {}\draw (\x,0.1)--(\x,-0.1) node[below]{$\x$};
			\foreach \y in {}\draw (0.1,\y)--(-0.1,\y) node[above left]{$\y$};
			\clip (\xmin,\ymin) rectangle (\xmax,\ymax);
			
			\draw[thick,smooth,samples=200,domain=\xmin:\xmax] plot (\x,{-1*(\x)+9});
			\clip (\xmin,\ymin) rectangle (\xmax,\ymax);
			
			\draw[thick,smooth,samples=200,domain=\xmin:\xmax] plot (\x,{(-1/4)*(\x)+6});
			
			\draw[thick,smooth,samples=200,domain=\xmin:\xmax] plot (\x,{-3*(\x)+21});
			
			\fill[pattern=north east lines,smooth]
			plot[domain=-3:10] (\x,{0*(\x)+10})
			-- plot[domain=10:-3] (\x,{-1*(\x)+9})
			-- cycle;
			
			\fill[pattern=north east lines,smooth]
			plot[domain=-3:10] (\x,{0*(\x)+10})
			-- plot[domain=10:-3] (\x,{(-1/4)*(\x)+6})
			-- cycle;
			
			\fill[pattern=north east lines,smooth]
			plot[domain=-3:10] (\x,{0*(\x)+10})
			-- plot[domain=10:-3] (\x,{-3*(\x)+21})
			-- cycle;
			
			\fill[pattern=north east lines,smooth]
			plot[domain=-3:10] (\x,{0*(\x)+0})
			-- plot[domain=10:-3] (\x,{0*(\x)-4})
			-- cycle;
			
			\fill[pattern=north east lines,smooth]
			plot[domain=-3:0] (\x,{0*(\x)+10})
			-- plot[domain=0:-3] (\x,{0*(\x)-2})
			-- cycle;
			
			\path
			(4,5) node[below left]{$A$}
			(6,3) node[below left]{$B$}
			(7,0) node[below left]{$C$}
			(0,0) node[below right]{$D$}
			(0,6) node[below right]{$E$}
			;
		\end{tikzpicture}
	\end{center}
		Miền nghiệm là ngũ giác $A B C D E$.\\
		Tọa độ các điểm: $A\left(4 ; 5\right)$, $B\left(6 ; 3\right)$, $C\left(7 ; 0\right)$, $D\left(0 ; 0\right)$, $E\left(0 ; 6\right)$.\\
		$M(x,y)$ sẽ đạt giá trị lớn nhất, giá trị nhỏ nhất tại các đỉnh của miền nghiệm nên thay tọa độ các điểm vào biểu thức $M\left(x, y\right)$ ta được\\
		$
		M\left(4 ; 5\right)=640$; $M\left(6 ; 3\right)=600$; $M\left(7 ; 0\right)=420$; $M\left(0 ; 0\right)=0$; $M\left(0 ; 6\right)=480$.\\
		
		Vậy giá trị lớn nhất của $M\left(x ; y\right)$ bằng $640$ khi $x=4$; $y=5 \Rightarrow a=4$; $b=5 \Rightarrow a-b=-1$.
		}
	\end{ex}
\Closesolutionfile{ans}
\indapan{6}{ans/ans-0-B1}
\cauds
\Opensolutionfile{ans}[ans/ans-0-B1-DS]

\begin{ex}%[0D2H2-1]
	Xét tính đúng, sai của các mệnh đề sau
	\choiceTF
	{$\heva{&x+\sqrt{y} > 3 \\ &x^{2}-y \ge -2}$ là hệ bất phương trình bậc nhất hai ẩn}
	{\True $\heva{&x > 5 \\ &y <-2 \\ & x+y \ge 100}$ là hệ bất phương trình bậc nhất hai ẩn}
	{\True $\heva{&3^{5}x-\sqrt{20}y > 7 \\ &x+y \le 100}$ là hệ bất phương trình bậc nhất hai ẩn}
	{$\heva{&x+y+z < 10 \\ &x+y < 5 \\& 2x+3y \ge 20}$ là hệ bất phương trình bậc nhất hai ẩn}
	\loigiai{
		\begin{itemchoice}
		\itemch Sai. Vì $\heva{&x+\sqrt{y} > 3 \\ &x^{2}-y \ge -2}$ chứa $\sqrt{x}$, $x^2$ nên không phải là hệ bất phương trình bậc nhất hai ẩn.
		\itemch Đúng. Vì $\heva{&x > 5 \\ &y <-2 \\ & x+y \ge 100}$ là hệ bất phương trình bậc nhất hai ẩn.
		\itemch Đúng. Vì $\heva{&3^{5}x-\sqrt{20}y > 7 \\ &x+y \le 100}$ là hệ bất phương trình bậc nhất hai ẩn
		\itemch Sai. Vì $\heva{&x+y+z < 10 \\ &x+y < 5 \\& 2x+3y \ge 20}$ chứa $z$ nên là hệ bất phương trình bậc nhất bai ẩn.
		\end{itemchoice}.}
\end{ex}

\begin{ex}%[0D2H2-2]
	Xét tính đúng, sai của các mệnh đề sau?
	\choiceTF
	{\True $\left(3;-1\right)$ là một nghiệm của hệ bất phương trình $\heva{&x+5y < 5 \\ &-3x-y \le 7}$}
	{\True $\left(3;-1\right)$ không là một nghiệm của hệ bất phương trình $\heva{&x > 5 \\ & y < 4 \\ &x+y \le 10}$}
	{\True $\left(3;-1\right)$ không là một nghiệm của hệ bất phương trình $\heva{&-x+5y > 1 \\ &3x+y > 5}$}
	{$\left(0;0\right)$ là một nghiệm của hệ bất phương trình $\heva{&2x+y > 3 \\ &-x+3y \le 5 \\ &3x-y \ge 7}$}
	\loigiai{
	\begin{itemchoice}
	\itemch Đúng. Vì thay $\left(3;-1\right)$ vào hệ bất phương trình $\heva{&x+5y < 5 \\ &-3x-y \le 7}$ ta được
	$\heva{&3+5\left(-1\right)<5 \\ &-3 \cdot 3 -\left(-1\right) \le -7}$ thoả mãn.
	\itemch Đúng. Vì thay $\left(3;-1\right)$ vào hệ bất phương trình $\heva{&x > 5 \\ & y < 4 \\ &x+y \le 10}$ ta được $\heva{&3 > 5 \\ & -1 < 4 \\ &3-1 \le 10}$ (vô lí).
	\itemch Đúng. Vì miền nghiệm của hệ bất phương trình $\heva{&-x+5y > 1 \\ &3x+y > 5}$ là miền không bị gạch ở hình sau (không kể bờ $d_{1}, d_{2}$)
\begin{center}
	\begin{tikzpicture}[scale=0.7, line join=round, line cap=round, >=stealth]
		\tikzset{every node/.style={scale=1}}
		\def\xmin{-3}\def\xmax{7}\def\ymin{-3}\def\ymax{7}
		\draw[->] (\xmin-0.2,0)--(\xmax+0.2,0) node[below]{$x$};
		\draw[->] (0,\ymin-0.2)--(0,\ymax+0.2) node[right]{$y$};
		\draw (0,0) node[below left]{$O$};
		\foreach \x in {-2,2,4,6}\draw (\x,0.1)--(\x,-0.1) node[below]{$\x$};
		\foreach \y in {-2,4,6}\draw (0.1,\y)--(-0.1,\y) node[above left]{$\y$};
		\clip (\xmin,\ymin) rectangle (\xmax,\ymax);
		
		\draw[thick,smooth,samples=200,domain=\xmin:\xmax] plot (\x,{-3*(\x)+5});
		\clip (\xmin,\ymin) rectangle (\xmax,\ymax);
		
		\draw[thick,smooth,samples=200,domain=\xmin:\xmax] plot (\x,{(1/5)*(\x)+1});
		
		
		\fill[pattern=north east lines,smooth]
		plot[domain=-3:4] (\x,{-3*(\x)+5})
		-- plot[domain=4:-3] (\x,{0*(\x)-10})
		-- cycle;
		
		\fill[pattern=north east lines,smooth]
		plot[domain=-3:11] (\x,{(1/5)*(\x)+1})
		-- plot[domain=11:-3] (\x,{0*(\x)-6})
		-- cycle;
		
	\end{tikzpicture}
\end{center}
	\itemch Sai. Vì miền nghiệm của hệ bất phương trình $\heva{&2x+y > 3 \\ &-x+3y \le 5 \\ &3x-y \ge 7}$ là miền không bị gạch ở hình sau (không kể bờ $d_{1}$).
	\begin{center}
		\begin{tikzpicture}[scale=0.7, line join=round, line cap=round, >=stealth]
			\tikzset{every node/.style={scale=1}}
			\def\xmin{-3}\def\xmax{7}\def\ymin{-3}\def\ymax{7}
			\draw[->] (\xmin-0.2,0)--(\xmax+0.2,0) node[below]{$x$};
			\draw[->] (0,\ymin-0.2)--(0,\ymax+0.2) node[right]{$y$};
			\draw (0,0) node[below left]{$O$};
			\foreach \x in {-2,2,4,6}\draw (\x,0.1)--(\x,-0.1) node[below]{$\x$};
			\foreach \y in {-2,4,6}\draw (0.1,\y)--(-0.1,\y) node[above left]{$\y$};
			\clip (\xmin,\ymin) rectangle (\xmax,\ymax);
			
			\draw[thick,smooth,samples=200,domain=\xmin:\xmax] plot (\x,{3*(\x)-7});
			\clip (\xmin,\ymin) rectangle (\xmax,\ymax);
			
			\draw[thick,smooth,samples=200,domain=\xmin:\xmax] plot (\x,{(1/3)*(\x)+(5/3)});
			
			\draw[thick,smooth,samples=200,domain=\xmin:\xmax] plot (\x,{-2*(\x)+3});
			
			\fill[pattern=north east lines,smooth]
			plot[domain=-5:6] (\x,{0*(\x)+7})
			-- plot[domain=6:-5] (\x,{3*(\x)-7})
			-- cycle;
			
			\fill[pattern=north east lines,smooth]
			plot[domain=-3:11] (\x,{0*(\x)+8})
			-- plot[domain=11:-3] (\x,{(1/3)*(\x)+(5/3)})
			-- cycle;
			
			\fill[pattern=north east lines,smooth]
			plot[domain=-3:11] (\x,{-2*(\x)+3})
			-- plot[domain=11:-3] (\x,{0*(\x)-5})
			-- cycle;
		\end{tikzpicture}
	\end{center}
\end{itemchoice}
	}
\end{ex}

\begin{ex}%[0D2V2-3]
	Bác Minh có kế hoạch đầu tư không quá $240$ triệu đồng vào hai khoản $X$ và khoản $Y$. Để đạt được lợi nhuận thì khoản $Y$ phài đầu tư it nhất 40 triệu đồng và số tiền đầu tư cho khoản $X$ phài it nhất gấp ba lần số tiền cho khoản $Y$. Khi đó:
	\choiceTF
	{\True Gọi $x$, $y$ (đơn vị: triệu đồng) tiền bác Minh đầu tư vào kho ta có hệ bất phương trình $\heva{&x+y \le 240 \\ &y \ge 40 \\ &x \ge 3y}$}
	{Miền nghiệm của hệ bất phương trình tiền bác Minh đầu tư vào kho là một tứ giác}
	{Điểm $C\left(200; 40\right)$ không thuộc miền nghiệm của hệ bất phương trình tiền bác Minh đầu tư vào kho}
	{\True Điểm $A\left(180; 60\right)$ là điểm có tung độ lớn nhất thuộc miền nghiệm của hệ bất phương trình tiền bác Minh đầu tư vào kho}
	\loigiai{
	\begin{itemchoice}
		
	
	\itemch Đúng. Vì gọi $x, y$ (đơn vị: triệu đồng) tiền bác Minh đầu tư vào kho Ta có hệ bất phương trình:$\heva{&x+y \le 240 \\ &y \ge 40 \\ &x \ge 3y}$.	
	\itemch Sai. Vì miền nghiệm của hệ trên là miền tam giác $A B C$ với $A(180 ; 60)$, $B(120 ; 40)$, $C(200 ; 40)$ ở Hình. 
	\begin{center}
		\begin{tikzpicture}[scale=0.03, line join=round, line cap=round, >=stealth]
			\tikzset{every node/.style={scale=1}}
			\def\xmin{-20}\def\xmax{300}\def\ymin{-20}\def\ymax{300}
			\draw[->] (\xmin-0.2,0)--(\xmax+0.2,0) node[below]{$x$};
			\draw[->] (0,\ymin-0.2)--(0,\ymax+0.2) node[right]{$y$};
			\draw (0,0) node[below left]{$O$};
			\foreach \x in {120,200,240}\draw (\x,0.1)--(\x,-0.1) node[below]{$\x$};
			\foreach \y in {40,80,240,280}\draw (0.1,\y)--(-0.1,\y) node[above left]{$\y$};
			\clip (\xmin,\ymin) rectangle (\xmax,\ymax);
			
			\draw[thick,smooth,samples=200,domain=\xmin:\xmax] plot (\x,{(1/3)*(\x)});
			\clip (\xmin,\ymin) rectangle (\xmax,\ymax);
			
			\draw[thick,smooth,samples=200,domain=\xmin:\xmax] plot (\x,{-1*(\x)+240});
			
			\draw[thick,smooth,samples=200,domain=\xmin:\xmax] plot (\x,{0*(\x)+40});
			
			\fill[pattern=north east lines,smooth]
			plot[domain=-40:300] (\x,{0*(\x)+300})
			-- plot[domain=300:-40] (\x,{(1/3)*(\x)})
			-- cycle;
			
			\fill[pattern=north east lines,smooth]
			plot[domain=-40:300] (\x,{0*(\x)+300})
			-- plot[domain=300:-40] (\x,{-1*(\x)+240})
			-- cycle;
			
			\fill[pattern=north east lines,smooth]
			plot[domain=-40:300] (\x,{0*(\x)+40})
			-- plot[domain=300:-40] (\x,{0*(\x)-40})
			-- cycle;
			
			\path
			(180,60) node[below left]{$A$}
			(120,40) node[below left]{$B$}
			(200,40) node[below left]{$C$}
			;
		\end{tikzpicture}
	\end{center}
	\itemch Sai. Vì điểm $C\left(200 ; 40\right)$ thuộc miền nghiệm của hệ bất phương trình tiền bác Minh đầu tư vào kho.
	\itemch Đúng. Vì điểm $A\left(180 ; 60\right)$ là điểm có tung độ lớn nhất thuộc miền nghiệm của hệ bất phương trình tiền bác Minh đầu tư vào kho
\end{itemchoice}
	}
\end{ex}

\begin{ex}%[0D2V2-2]
	Cho hệ bất phương trình: $\heva{&3x+2y \ge 9 \\ &x-2y \le 3 \\ &x+y \le 6 \\ &x \ge 1}$ (I). Khi đó:
	\choiceTF
	{Miền nghiệm của hệ bất phương trình là miền tam giác}
	{\True $\left(3; 2\right)$ là một nghiệm của hệ bất phương trình}
	{$x=1, y=3$ là nghiệm của hệ bất phương trình (I) sao cho $F=3 x-y$ đạt giá trị lớn nhất}
	{\True $x=1, y=5$ là nghiệm của hệ bất phương trình (I) sao cho $F=3 x-y$ đạt giá trị nhỏ nhất}
	\loigiai{
	\begin{itemchoice}
		\itemch Sai. Vì miền nghiệm của hệ (I) là miền tứ giác $A B C D$ với $A\left(3 ; 0\right)$, $B\left(5 ; 1\right)$, $C\left(1 ; 5\right)$, $D\left(1 ; 3\right)$ (Hình).
	\begin{center}
		\begin{tikzpicture}[scale=0.7, line join=round, line cap=round, >=stealth]
			\tikzset{every node/.style={scale=1}}
			\def\xmin{-3}\def\xmax{7}\def\ymin{-3}\def\ymax{8}
			\draw[->] (\xmin-0.2,0)--(\xmax+0.2,0) node[below]{$x$};
			\draw[->] (0,\ymin-0.2)--(0,\ymax+0.2) node[right]{$y$};
			\draw (0,0) node[below left]{$O$};
			\foreach \x in {1,3,5,6}\draw (\x,0.1)--(\x,-0.1) node[below]{$\x$};
			\foreach \y in {1,3,5,6}\draw (0.1,\y)--(-0.1,\y) node[above left]{$\y$};
			\clip (\xmin,\ymin) rectangle (\xmax,\ymax);
			
			\draw[thick,smooth,samples=200,domain=\xmin:\xmax] plot (\x,{(-3/2)*(\x)+(9/2)});
			\clip (\xmin,\ymin) rectangle (\xmax,\ymax);
			
			\draw[thick,smooth,samples=200,domain=\xmin:\xmax] plot (\x,{(1/2)*(\x)-(3/2)});
			
			\draw[thick,smooth,samples=200,domain=\xmin:\xmax] plot (\x,{-1*(\x)+6});
			
			\draw[thick,smooth,samples=200,domain=\xmin:\xmax] (1,-3) -- (1,8);
			
			\fill[pattern=north east lines,smooth]
			plot[domain=-3:7] (\x,{(-3/2)*(\x)+(9/2)})
			-- plot[domain=7:-3] (\x,{0*(\x)-5})
			-- cycle;
			
			\fill[pattern=north east lines,smooth]
			plot[domain=-3:7] (\x,{(1/2)*(\x)-(3/2)})
			-- plot[domain=7:-3] (\x,{0*(\x)-5})
			-- cycle;
			
			\fill[pattern=north east lines,smooth]
			plot[domain=-3:7] (\x,{0*(\x)+8})
			-- plot[domain=7:-3] (\x,{-1*(\x)+6})
			-- cycle;
			
			\fill[pattern=north east lines,smooth]
			plot[domain=-3:1] (\x,{0*(\x)+8})
			-- plot[domain=1:-3] (\x,{0*(\x)-3})
			-- cycle;
			
			\path
			(3,0) node[above]{$A$}
			(5,1) node[above]{$B$}
			(1,5) node[above right]{$C$}
			(1,3) node[above right]{$D$}
			;
		\end{tikzpicture}
	\end{center}
		\itemch Đúng. Vì $\left(3 ; 2\right)$ là một nghiệm của hệ bất phương trình.
		\itemch Sai. Vì tính giá trị của $F=3 x-y$ tại các cặp số $(x ; y)$ là toạ độ của các đỉnh tứ giác $A B C D$ rồi so sánh các giá trị đó, ta được $F$ đạt giá trị lớn nhất bằng 14 tại $x=5$, $y=1$
		\itemch Đúng. Vì $F$ đạt giá trị nhỏ nhất bằng $-2$ tại $x=1$, $y=5$.
	\end{itemchoice}
	}
\end{ex}

\Closesolutionfile{ans}
\indapan{3}{ans/ans-0-B1-DS}

\Opensolutionfile{ans}[ans/ans-0-B1-KQ]
\caukq

\begin{ex}%[0D2H2-2]
	Cho hệ bất phương trình: $\heva{&x+y \ge 5 \\ &x-2y \le 2 \\ &y \le 3}$ (II). Giá trị của tham số $m \ge a$ đề bất phương trình $2 x-5 y+m \geq 0$  nghiệm đúng với mọi cặp số $\left(x; y\right)$ thoả mãn hệ bất phương trình (II). Tìm $a$
	\shortans{$11$}
	\loigiai{Miền nghiệm của hệ bất phương trình (II) là miền tam giác $A B C$ với $A\left(4 ; 1\right), B\left(8 ; 3\right), C\left(2 ; 3\right)$ (Hình).
\begin{center}
	\begin{tikzpicture}[scale=0.7, line join=round, line cap=round, >=stealth]
		\tikzset{every node/.style={scale=1}}
		\def\xmin{-1}\def\xmax{10}\def\ymin{-2}\def\ymax{8}
		\draw[->] (\xmin-0.2,0)--(\xmax+0.2,0) node[below]{$x$};
		\draw[->] (0,\ymin-0.2)--(0,\ymax+0.2) node[right]{$y$};
		\draw (0,0) node[below right]{$O$};
		\foreach \x in {2,4,5,8}\draw (\x,0.1)--(\x,-0.1) node[below]{$\x$};
		\foreach \y in {-1,3,5}\draw (0.1,\y)--(-0.1,\y) node[left]{$\y$};
		\clip (\xmin,\ymin) rectangle (\xmax,\ymax);
		
		\draw[thick,smooth,samples=200,domain=\xmin:\xmax] plot (\x,{-1*(\x)+5});
		
		\draw[thick,smooth,samples=200,domain=\xmin:\xmax] plot (\x,{(1/2)*(\x)-1});
		
		\draw[thick,smooth,samples=200,domain=\xmin:\xmax] plot (\x,{0*(\x)+3});
		
		%\draw[thick,smooth,samples=200,domain=\xmin:\xmax] (1.6,-2) -- (1.6,3);
		
		\fill[pattern=north east lines,smooth]
		plot[domain=-1:10] (\x,{-1*(\x)+5})
		-- plot[domain=10:-1] (\x,{0*(\x)-2})
		-- cycle;
		
		\fill[pattern=north east lines,smooth]
		plot[domain=-1:10] (\x,{(1/2)*(\x)-1})
		-- plot[domain=10:-1] (\x,{0*(\x)-2})
		-- cycle;
		
		
		\fill[pattern=north east lines,smooth]
		plot[domain=-1:10] (\x,{0*(\x)+10})
		-- plot[domain=10:-1] (\x,{0*(\x)+3})
		-- cycle;
		
		
		
		\path
		(4,1) node[above]{$A$}
		(8,3) node[below]{$B$}
		(2,3) node[below]{$C$}
		;
	\end{tikzpicture}		
\end{center}
	Ta có: $2 x-5 y+m \ge 0 \Leftrightarrow m \ge-2 x+5 y$.\\		
	Đặt $F=-2 x+5 y$. Tính giá trị của $F=-2 x+5 y$ tại các cặp số $(x ; y)$ là toạ độ của các đĩnh tam giác $A B C$ rồi so sánh các giá trị đó, ta được $F$ đạt giá trị lớn nhất bằng 11 tại $x=2, y=3$.\\		
	Để bất phương trình $2 x-5 y+m \geq 0$ nghiệm đúng với mọi $x, y$ thoả mãn hệ bất phương trình đã cho thì $m \ge \max  F$ trên miền nghiệm của hệ bất phương trình đó hay $m \ge 11$.\\	
	Vậy $a=11$.}
	\end{ex}

	\begin{ex}%[0D2V2-3]
		Trong một cuộc thi pha chế đồ uống gồm hai loại là $A$ và $B$, mỗi đội chơi được sử dụng tối đa $24$ g hương liệu, $9$ cốc nước lọc và $210$ g đường. Để pha chế $1$ cốc đồ uống loại $A$ cần $1$ cốc nước lọc, $30$ g đường và $1$ g hương liệu. Để pha chế 1 cốc đồ uống loại $B$ cần 1 cốc nước lọc, $10$ g đường và $4$ g hương liệu. Mỗi cốc đồ uống loại $A$ nhận được 6 điểm thương, mỗi cốc đồ uống loại $B$ nhận được 8 điểm thường. Để đạt được số điểm thưởng cao nhất, đội chơi cần pha chế $ x $ cốc đồ uống loại $A, y $ cốc đồ uống loại $B$. Tính $ x+y $ 
		\shortans{$9$}
	\loigiai{
		Gọi $x, y$ lần lượt là số cốc đồ uống loại $A$, loại $B$ mà đội chơi cần pha chế với $x \ge 0, y \ge 0$.\\	
		Số cốc nước cần dùng là $x+y$ (cốc).\\	
		Lượng đường cần dùng là $30 x+10 y$ (g).\\	
		Lượng hương liệu cần dùng là $x+4 y$ (g).\\	
		Theo giả thiết, ta có $\heva{&x \ge 0 \\ &y \ge 0 \\ &x+y \le 9 \\ &30x+10y \le 210 \\ &x+4y \le 24} \Leftrightarrow \heva{&x \ge 0 \\ &y \ge 0 \\ &x+y \le 9 \\ &3x+y \le 21 \\ &x+4y \le 24}$ (III)\\	
		Số điểm thường nhận được là: $F=6 x+8 y$.\\	
		Ta tìm giá trị lớn nhất trên miền nghiệm của hệ bất phương trình (III).\\	
		Miền nghiệm của hệ bất phương trình (III) là miền ngũ giác $O A B C D$ với $O\left(0 ; 0\right)$, $A\left(7 ; 0\right)$, $B\left(6 ; 3\right)$, $C\left(4 ; 5\right)$, $D\left(0 ; 6\right)$ (hình).
	\begin{center}
		\begin{tikzpicture}[scale=0.7, line join=round, line cap=round, >=stealth]
			\tikzset{every node/.style={scale=1}}
			\def\xmin{-1}\def\xmax{10}\def\ymin{-2}\def\ymax{10}
			\draw[->] (\xmin-0.2,0)--(\xmax+0.2,0) node[below]{$x$};
			\draw[->] (0,\ymin-0.2)--(0,\ymax+0.2) node[right]{$y$};
			\draw (0,0) node[below right]{$O$};
			
			\clip (\xmin,\ymin) rectangle (\xmax,\ymax);
			
			\draw[thick,smooth,samples=200,domain=\xmin:\xmax] plot (\x,{-1*(\x)+9});
			
			\draw[thick,smooth,samples=200,domain=\xmin:\xmax] plot (\x,{-3*(\x)+21});
			
			\draw[thick,smooth,samples=200,domain=\xmin:\xmax] plot (\x,{(-1/4)*(\x)+6});
			
			%\draw[thick,smooth,samples=200,domain=\xmin:\xmax] (1.6,-2) -- (1.6,3);
			
			\fill[pattern=north east lines,smooth]
			plot[domain=-1:10] (\x,{0*(\x)+10})
			-- plot[domain=10:-1] (\x,{-1*(\x)+9})
			-- cycle;
			
			\fill[pattern=north east lines,smooth]
			plot[domain=4:10] (\x,{0*(\x)+10})
			-- plot[domain=10:4] (\x,{-3*(\x)+21})
			-- cycle;
			
			
			\fill[pattern=north east lines,smooth]
			plot[domain=-1:10] (\x,{0*(\x)+10})
			-- plot[domain=10:-1] (\x,{(-1/4)*(\x)+6})
			-- cycle;
			
			\fill[pattern=north east lines,smooth]
			plot[domain=-1:0] (\x,{0*(\x)+10})
			-- plot[domain=0:-1] (\x,{0*(\x)-2})
			-- cycle;
			
			\fill[pattern=north east lines,smooth]
			plot[domain=-1:10] (\x,{0*(\x)+0})
			-- plot[domain=10:-1] (\x,{0*(\x)-2})
			-- cycle;

			\path
			(7,0) node[above left]{$A$}
			(6,3) node[below left]{$B$}
			(4,5) node[below left]{$C$}
			(0,6) node[below right]{$D$}
			;
			\draw[dashed] (6,0) node[below]{$6$}
			-- (6,3)  --
			(0,3) node[left]{$3$};
			\draw[dashed] (4,0) node[below]{$4$}
			-- (4,5)  --
			(0,5) node[left]{$5$};
		\end{tikzpicture}		
	\end{center}
		Tính giá trị của $F=6 x+8 y$ tại các cặp số $\left(x, y\right)$ là tọa độ của các đỉnh ngũ giác $O A B C D$ rồi so sánh các giá trị đó, ta được $F$ đạt giá trị lớn nhất bằng $64$ tại $x=4 ; y=5$.\\		
		Để đạt được số điểm thưởng cao nhất, đội chơi cần pha chế 4 cốc đồ uống loại $A, 5$ cốc đồ uống loại $B$.\\		
		Vậy $x+y=4+5=9$.
	}
	\end{ex}
	
	\begin{ex}%[0D2V2-3]
		Một gia đình cần ít nhất $900$ đơn vị protein và $400$ đơn vị lipit trong thức ăn mỗi ngày. Mỗi ki-lô-gam thịt bò chứa $800$ đơn vị protein và $200$ đơn vị lipit. Mỗi ki-lô-gam thịt lợn (heo) chứa $600$ đơn vị protein và $400$ đơn vị lipit. Biết rằng gia đình này chỉ mua nhiều nhất là $1{,}6$ kg thịt bò và $1{,}1$ kg thịt lợn; giá $1$ kg thịt bò là $200\,000$ đồng, $1$ kg thịt lợn là $160\,000$ đồng. Gia đình đó cần mua $x$ kg thịt bò và $y$ kg thịt lợn để đảm bảo cung cấp đủ lượng protein, lipit cho gia đình và có chi phí là ít nhất. Tính $ a+b $.
		\shortans{$1{,}3$}
		
	\loigiai{
		Gọi $x, y$ lần lượt là số ki-lô-gam thịt bò và thịt lợn mà gia đình đó mua trong một ngày với $0 \le x \le 1{,}6, 0 \le y \le 1{,}1$ Số đơn vị protein gia đình có là $800 x+600 y$.\\		
		Số đơn vị lipit gia đình có là $200 x+400 y$. Theo bài ra, ta có
		$$\heva{&0 \le x \le 1{,}6 \\ &0 \le y \le 1{,}1 \\ &800x+600y \ge 900 \\ &200x+400y \ge 400} \Leftrightarrow \heva{&0 \le x \le 1{,}6 \\ &0 \le y \le 1{,}1 \\ &8x+6y \ge 9 \\ &x+2y \ge 2}$$ 
		Số tiền gia đình đã dùng để mua thịt bò và thịt lợn là\\
		$
		T=200\,000 x+160\,000 y
		$ (đồng)

		Bài toán đưa về tìm $x, y$ là nghiệm của hệ bất phương trình  để $T=200\,000 x+160\,000 y$ đạt giá trị nhỏ nhất.\\		
		Trước hết, ta biểu diễn miền nghiệm của hệ bất phương trình.\\
		Miền nghiệm của hệ bất phương trình là miền tứ giác $ABCD$ với $A\left(0{,}3;1{,}1\right)$, $B\left(0{,}6;0{,}7\right)$, $C\left(1{,}6;0{,}2\right)$, $D\left(1{,}6;1{,}1\right)$ (hình).
	\begin{center}
		\begin{tikzpicture}[scale=1.5, line join=round, line cap=round, >=stealth]
			\tikzset{every node/.style={scale=1}}
			\def\xmin{-1}\def\xmax{3}\def\ymin{-2}\def\ymax{3}
			\draw[->] (\xmin-0.2,0)--(\xmax+0.2,0) node[below]{$x$};
			\draw[->] (0,\ymin-0.2)--(0,\ymax+0.2) node[right]{$y$};
			\draw (0,0) node[below right]{$O$};
			
			\clip (\xmin,\ymin) rectangle (\xmax,\ymax);
			
			\draw[thick,smooth,samples=200,domain=\xmin:\xmax] plot (\x,{0*(\x)+1.1});
			\draw[thick,smooth,samples=200,domain=\xmin:\xmax] plot (\x,{(-4/3)*(\x)+3/2});
			\draw[thick,smooth,samples=200,domain=\xmin:\xmax] plot (\x,{(-1/2)*(\x)+1});
			
			\draw[thick,smooth,samples=200,domain=\xmin:\xmax] (1.6,-2) -- (1.6,3);
			
			\fill[pattern=north east lines,smooth]
			plot[domain=-2:3] (\x,{0*(\x)+3})
			-- plot[domain=3:-2] (\x,{0*(\x)+1.1})
			-- cycle;
			
			\fill[pattern=north east lines,smooth]
			plot[domain=-1:3] (\x,{(-4/3)*(\x)+3/2})
			-- plot[domain=3:-1] (\x,{0*(\x)-2})
			-- cycle;
			
			
			\fill[pattern=north east lines,smooth]
			plot[domain=-1:3] (\x,{(-1/2)*(\x)+1})
			-- plot[domain=3:-1] (\x,{0*(\x)-2})
			-- cycle;
			
			\fill[pattern=north east lines,smooth]
			plot[domain=1.6:3] (\x,{0*(\x)+3})
			-- plot[domain=3:1.6] (\x,{0*(\x)-2})
			-- cycle;
			
			\draw[dashed] (0,1.1)
			 -- (0.3,1.1) node[above right]{$A$} --
			(0.3,0);
			
			\draw[dashed] (0,0.7)
			-- (0.6,0.7) node[above right]{$B$} --
			(0.6,0);
			\path
			(1.6,1.1) node[above right]{$D$}
			(1.6,0.2) node[above right]{$C$}
			;
		\end{tikzpicture}		
	\end{center}
		Tính giá trị của $T$ tại các cặp sổ $\left(x, y\right)$ là tọa độ của các đỉnh tứ giác $A B C D$ rồi so sánh các giá trị đó, ta được $T$ đạt giá trị nhỏ nhất bằng $232\,000$ đồng tại $x=0{,}6 ; y=0{,}7$.\\		
		Để đảm bảo cung cấp đủ lượng protein, lipit cho gia đình và có chi phí là ít nhất thì gia đình đó cần mua thêm $0{,}6$ kg thịt bò và $0{,}7$ kg thịt lợn.\\		
		Vậy $x+y=0{,}6+0{,}7=1{,}3$.
	}
	\end{ex}
	
	\begin{ex}%[0D2V2-1]
		Tìm GTLN của $f\left(x, y\right)=x+2 y$ với điều kiện 
		$$\heva{&0 \le y \le 4 \quad\quad\quad  \left(d_{1}\right)\\&0 \le x \quad\quad\quad\quad\quad \left(d_{2}\right)\\ &x-y-1 \le 0 \quad\quad \left(d_{3}\right)\\ &x+2y-10 \le 0 \quad \left(d_{4}\right)}$$
		\shortans{$10$}
	\loigiai{
	\begin{center}
		\begin{tikzpicture}[scale=0.7, line join=round, line cap=round, >=stealth]
			\tikzset{every node/.style={scale=1}}
			\def\xmin{-3}\def\xmax{7}\def\ymin{-1}\def\ymax{6}
			\draw[->] (\xmin-0.2,0)--(\xmax+0.2,0) node[below]{$x$};
			\draw[->] (0,\ymin-0.2)--(0,\ymax+0.2) node[right]{$y$};
			\draw (0,0) node[below right]{$O$};
			
			\clip (\xmin,\ymin) rectangle (\xmax,\ymax);
			
			\draw[thick,smooth,samples=200,domain=\xmin:\xmax] plot (\x,{0*(\x)+4});
			\draw[thick,smooth,samples=200,domain=\xmin:\xmax] plot (\x,{1*(\x)-1});
			\draw[thick,smooth,samples=200,domain=\xmin:\xmax] plot (\x,{(-1/2)*(\x)+5});
			
			
			\fill[pattern=north east lines,smooth]
			plot[domain=-3:7] (\x,{0*(\x)+6})
			-- plot[domain=7:-3] (\x,{0*(\x)+4})
			-- cycle;
			
			\fill[pattern=north east lines,smooth]
			plot[domain=-2:7] (\x,{0*(\x)+6})
			-- plot[domain=7:-2] (\x,{(-1/2)*(\x)+5})
			-- cycle;
			
			\fill[pattern=north east lines,smooth]
			plot[domain=0:7] (\x,{1*(\x)-1})
			-- plot[domain=7:0] (\x,{0*(\x)-2})
			-- cycle;
			
			\fill[pattern=north east lines,smooth]
			plot[domain=-5:0] (\x,{0*(\x)+6})
			-- plot[domain=0:-5] (\x,{0*(\x)-2})
			-- cycle;
			
			\fill[pattern=north east lines,smooth]
			plot[domain=-5:7] (\x,{0*(\x)})
			-- plot[domain=7:-5] (\x,{0*(\x)-2})
			-- cycle;
			
			\path
			(0,4) node[below right]{$D$}
			(1,0) node[above left]{$A$}
			(4,3) node[left]{$B$}
			(2,4) node[below]{$C$}
			;
		\end{tikzpicture}		
	\end{center}
		$M\left(x; y\right)$ thoả mãn (I) là miền bên trong đa giác $OABCD$.\\		
		Tìm toạ độ $A, B, C, D$ bằng phương pháp đồ thị hay phương trình hoành độ giao điểm.\\		
		Thay toạ độ $O, A, B, C, D$ vào $f\left(x; y\right)=x+2 y$ ta có
		\begin{center}
			\begin{tabular}{|c|c|c|c|c|c|}
				\hline & 0 & $A$ & $B$ & $C$ & $D$ \\
				\hline$f(x ; y)$ & $0$ & $1$ & $10$ & $10$ & $8$ \\
				\hline
			\end{tabular}
		\end{center}
		Vậy $\max f\left(x\right)=10$.
	}
	\end{ex}
	
	\begin{ex}%[0D2V2-2]
		Cho biểu thức $T=3 x-2 y-4$ với $x$ và $y$ thỏa mãn hệ bất phương trình $\heva{&x-y-1 \le 0 \\ &x+4y+9 \ge 0 \\ &x-2y+3 \ge 0}$.
		Biết $T$ đạt giá trị nhỏ nhất khi $x=x_0$ và $y=y_0$. Tính $x_0^2+y_0^2$.
		\shortans{$26$}
	\loigiai{
		Miền nghiệm của hệ là miền tam giác $ABC$ với $A\left(-5;-1\right), B\left(-1;-2\right)$ và $C \left(5; 4\right)$. Lập bảng		
		\begin{center}
			\begin{tabular}{|c|c|c|c|}
				\hline Đỉnh & $A (-5;-1)$ & $B(-1;-2)$ & $C (5; 4)$ \\
				\hline$T$ &$-17$ &$-3$ & $3$ \\
				\hline
			\end{tabular}
		\end{center}
		\begin{center}
		\begin{tikzpicture}[scale=0.7, line join=round, line cap=round, >=stealth]
			\tikzset{every node/.style={scale=1}}
			\def\xmin{-6}\def\xmax{6}\def\ymin{-3}\def\ymax{5}
			\draw[->] (\xmin-0.2,0)--(\xmax+0.2,0) node[below]{$x$};
			\draw[->] (0,\ymin-0.2)--(0,\ymax+0.2) node[right]{$y$};
			\draw (0,0) node[above right]{$O$};
			
			\clip (\xmin,\ymin) rectangle (\xmax,\ymax);
			\draw[thick,smooth,samples=200,domain=\xmin:\xmax] plot (\x,{1*(\x)-1});
			\draw[thick,smooth,samples=200,domain=\xmin:\xmax] plot (\x,{(-1/4)*(\x)-(9/4)});
			\draw[thick,smooth,samples=200,domain=\xmin:\xmax] plot (\x,{(1/2)*(\x)+(3/2)});
			
			\draw[dashed] (0,4)
			node[left]{$4$} -- (5,4) node[above left]{$C$} --
			(5,0) node[below]{$5$};
			
			\draw[dashed] (0,-1)
			node[right]{$-1$} -- (-5,-1) node[below]{$A$} --
			(-5,0) node[above]{$-5$};
			
			\draw[dashed] (0,-2)
			node[right]{$-2$} -- (-1,-2) node[below]{$B$} --
			(-1,0) node[above]{$-1$};
			
			\fill[pattern=north east lines,smooth]
			plot[domain=-1:6] (\x,{1*(\x)-1})
			-- plot[domain=6:-1] (\x,{(-1/4)*(\x)-(9/4)})
			-- cycle;
			\fill[pattern=north east lines,smooth]
			plot[domain=-6:3] (\x,{(-1/4)*(\x)-(9/4)})
			-- plot[domain=3:-6] (\x,{0*(\x)-3})
			-- cycle;
			\fill[pattern=north east lines,smooth]
			plot[domain=-6:6] (\x,{0*(\x)+5})
			-- plot[domain=6:-6] (\x,{(1/2)*(\x)+(3/2)})
			-- cycle;
		\end{tikzpicture}
		\end{center}
		Vậy $T$ đạt giá trị nhỏ nhất bằng $-17$ khi $x=-5$ và $y=-1$.\\
		Do đó $x_0=-5$ và $y_0=-1 \Rightarrow x_0^2+y_0^2=26$.
	}
	\end{ex}
	\begin{ex}%[0D2V2-3]
		 \immini{Trong mặt phẳng $O x y$, cho tứ giác $A B C D$ có $A(-3 ; 0)$; $B(0 ; 2)$; $C(3 ; 1)$; $D(3 ;-2)$. Tìm tất cả các giá trị nguyên của $m$ sao cho điểm $M(m ; m-1)$ nằm trong hình tứ giác $A B C D$ kể cả $4$ cạnh.}
		{\begin{tikzpicture}[scale=1, line join=round, line cap=round, >=stealth]
					\tikzset{every node/.style={scale=1}}
					\def\xmin{-3.5}\def\xmax{4}\def\ymin{-3.5}\def\ymax{4}
					\draw[->] (\xmin-0.2,0)--(\xmax+0.2,0) node[below]{$x$};
					\draw[->] (0,\ymin-0.2)--(0,\ymax+0.2) node[right]{$y$};
					\draw (0,0) node[below right]{$O$};
					
					\clip (\xmin,\ymin) rectangle (\xmax,\ymax);
					
					\draw[line width=1.3,fill=black] (-3,0)circle(2pt) node[above]{$A$}--(0,2)circle(2pt)node[above left]{$B$}--(3,1)circle(2pt)node[above]{$C$}--(3,-2)circle(2pt)node[below]{$D$}--(-3,0)(1.3,0.3)circle(2pt)node[right]{$M$};
				\foreach \a in {-2,-1,1,2} \fill (\a,0) circle (1.2pt);
				\foreach \b in {-3,-2,1,3} \fill (0,\b) circle (1.2pt);
				\end{tikzpicture}		
		}
		\shortans{$3$}
	\loigiai{
	Nhận thấy hình tứ giác $A B C D$ tính cả $4$ cạnh của nó là miền nghiệm của hệ bất phương trình gồm $4$ bất phương trình có miền nghiệm là nửa mặt phẳng chứa điểm $O(0 ; 0)$ và lần lượt có các bờ là các đường thẳng $A B$, $B C$, $C D$ và $D A$.\\	
	Phương trình đường thẳng $A B$ 
	$$
	\dfrac{x+3}{0-(-3)}=\dfrac{y-0}{2-0} \Leftrightarrow 2 x-3 y+6=0.
	$$
	
	Bất phương trình có miền nghiệm là là nửa mặt phẳng bờ $A B$ (tính cả bờ $A B$) và chứa điểm $O$ là $2 x-3 y+6 \geq 0$.\\	
	Phương trình đường thẳng $B C\colon  \dfrac{x-0}{3-0}=\dfrac{y-2}{1-2} \Leftrightarrow x+3 y-6=0$.\\ Bất phương trình có miền nghiệm là là nửa mặt phẳng bờ $B C$ (tính cả bờ $B C$) và chứa điểm $O$ là $x+3 y-6 \leq 0$.\\
	Phương trình đường thẳng $C D\colon  x-3=0$. Bất phương trình có miền nghiệm là là nửa mặt phẳng bờ $C D$ (tính cả bờ $C D$) và chứa điểm $O$ là $x-3 \leq 0$.\\	
	Phương trình đường thẳng $D A\colon  \dfrac{x+3}{3-(-3)}=\dfrac{y-0}{-2-0} \Leftrightarrow x+3 y+3=0$.\\
	 Bất phương trình có miền nghiệm là là nửa mặt phẳng bờ $D A$ (tính cả bờ $D A$) và chứa điểm $O$ là $x+3 y+3 \geq 0$.\\	
	Hình tứ giác $A B C D$ tính cả $4$ cạnh của nó là miền nghiệm của hệ bất phương trình $\heva{&2 x-3 y+6 \geq 0 \\& x+3 y-6 \leq 0 \\& x-3 \leq 0 \\& x+3 y+3 \geq 0.}$\, \hfill{(1)}\\
	Điểm $M(m ; m-1)$ nằm trong hình tứ giác $A B C D$ tính cả $4$ cạnh của nó khi và chỉ khi $(m ; m-1)$ là một nghiệm của hệ $(1)$, tức là
	$$
	\heva{& 2 m - 3 ( m - 1 ) + 6 \geq 0\\& m + 3 ( m - 1 ) - 6 \leq 0 \\& m - 3 \leq 0 \\& m + 3 ( m - 1 ) + 3 \geq 0 }
	\Leftrightarrow \heva{&	m \leq 9 \\&	m \leq \dfrac{9}{4} \\&	m \leq 3 \\&m \geq 0
	} \Leftrightarrow 0 \leq m \leq \dfrac{9}{4}.
	$$
	Vậy các giá trị của $m$ thỏa mãn là $0 \leq m \leq \dfrac{9}{4}$.
	}
	\end{ex}

\Closesolutionfile{ans}
\indapan{6}{ans/ans-0-B1-KQ}
\newpage 
\subsection{Đề 2}
\Opensolutionfile{ans}[ans/ans-0-De2]
\caulc
\begin{ex}%[0D3N1-2]
	Cặp số nào sau đây là nghiệm của hệ bất phương trình  $\heva{& x+2y\le8\\ & 3x-y>3 }$ ?
	\choice
	{$\left(0; 4\right)$ }
	{$\left(1; -1\right)$}
	{\True $\left(4; 1\right)$}
	{$\left(2; 3\right)$}
	\loigiai{
		Lần lượt thay các bộ số vào hệ bất phương trình ta được một nghiệm của hệ bất phương trình trên là $\left(4; 1\right)$.
	}
\end{ex}	

\begin{ex}%[0D3N1-2]
	Cho hệ bất phương trình $\heva{& x+3y-2\ge0\\ & 2x+y+1\le0 }$. Trong các điểm sau, điểm nào thuộc miền nghiệm của hệ bất phương trình?
	\choice
	{$M\left(0;1\right)$}
	{\True $N\left(-1;1\right)$}
	{$P\left(1;3\right)$}
	{$Q\left(-1;0\right)$}
	\loigiai{
		Ta thay lần lượt tọa độ các điểm vào hệ bất phương trình.\\
		Ta thấy $N\left(-1;1\right)$ thỏa mãn.
	}
\end{ex}

\begin{ex}%[0D3N1-2]
	Miền nghiệm của hệ bất phương trình $\heva{& x-2y+2>0\\ & 2x+y>3 }$  là phần mặt phẳng chứa điểm
	\choice
	{$\left(1;1\right)$ }
	{$\left(-1;2\right)$}
	{$\left(2;-1\right)$}
	{\True $\left(1;2\right)$}
	\loigiai{
		Lần lượt thay các bộ số vào hệ bất phương trình ta được một nghiệm của hệ bất phương trình trên là  $\left(1;2\right)$.
	}
\end{ex}	

\begin{ex}%[0D2H2-2]
	\immini[thm]{
		Phần \textbf{không tô đậm} trong hình vẽ bên (không chứa biên), biểu diễn tập nghiệm của hệ bất phương trình nào trong các hệ bất phương trình sau?
		\choice
		{$\heva{& x-y\ge 0 \\& 2x-y\ge 1}$}
		{\True $\heva{& x-y>0 \\& 2x-y>1}$}
		{$\heva{& x-y<0 \\& 2x-y>1}$}
		{$\heva{& x-y<0 \\& 2x-y<1}$}
	}
	{
		\begin{tikzpicture}[line join=round, line cap=round, >=stealth,font=\footnotesize, scale=0.7]
			\draw[->](-2,0)--(3,0) node[below right] {$x$};
			\draw[->](0,-3)--(0,3) node[above] {$y$};
			\clip (-2,-3) rectangle (3,3);
			\node (0,0) [above left]{$ O $};
			\foreach \x in {-1,...,2}
			\draw[shift={(\x,0)},color=black] (0pt,2pt) -- (0pt,-2pt);
			\foreach \y in {-1,...,3}
			\draw[shift={(0,\y)},color=black] (2pt,0pt) -- (-2pt,0pt);
			\draw[samples=100,smooth,domain=-4:2] plot(\x,{(\x)});
			\draw[samples=100,smooth,domain=-4:2] plot(\x,{2*(\x)-1});
			\draw [pattern=north west lines] (3,3)--(-2,3)--(-2,-3)--(-1,-3)--(1,1)--cycle;
			\draw[dashed] (0,1)--(1,1)--(1,0);
			\draw (0,1) node[left]{$1$} (1,0) node[below]{$1$} (0,-1) node[left]{$-1$};
		\end{tikzpicture}
	}
	\loigiai
	{
		Do miền nghiệm không chứa biên nên ta loại đáp án $\heva{& x-y\ge 0 \\& 2x-y\ge 1.}$\\
		Chọn điểm $M(1;0)$ thử vào các hệ bất phương trình.\\
		Xét đáp án $\heva{& x-y>0 \\& 2x-y>1}$, ta có $\heva{& 1-0>0 \\& 2\cdot 1-0>1}$ (luôn đúng) và miền nghiệm không chứa biên.
	}
\end{ex}


\begin{ex}%[0D2H2-2]
	\immini[thm]{
		Phần \textbf{không tô đậm} trong hình vẽ bên (không chứa biên), biểu diễn tập nghiệm của hệ bất phương trình nào trong các hệ bất phương trình sau?
		\choice
		{$\heva{& x-2y\le 0 \\& x+3y\ge -2}$}
		{$\heva{& x-2y>0 \\& x+3y<-2}$}
		{$\heva{& x-2y\le 0 \\& x+3y\le -2}$}
		{\True $\heva{& x-2y<0 \\& x+3y>-2}$}
	}
	{
		\begin{tikzpicture}[line join=round, line cap=round, >=stealth,font=\footnotesize, scale=0.8]
			\draw[->](-3,0)--(3,0) node[below right] {$x$};
			\draw[->](0,-2)--(0,2) node[above] {$y$};
			\clip (-3,-2) rectangle (3,2);
			\node (0,0) [above left]{$ O $};
			\foreach \x in {-2,...,2}
			\draw[shift={(\x,0)},color=black] (0pt,2pt) -- (0pt,-2pt);
			\foreach \y in {-1,...,1}
			\draw[shift={(0,\y)},color=black] (2pt,0pt) -- (-2pt,0pt);
			\draw[samples=100,smooth,domain=-3:3] plot(\x,{(\x)/2});
			\draw[samples=100,smooth,domain=-3:3] plot(\x,{(-(\x)-2)/3});
			\draw [pattern=north east lines] (-3,1/3)--(-3,-2)--(3,-2)--(3,3/2)--(-0.8,-0.4)--cycle;
			\draw[dashed] (0,1)--(2,1)--(2,0);
			\draw (0,1) node[left]{$1$} (2,0) node[below]{$2$} (-2,0) node[above]{$-2$};
		\end{tikzpicture}
	}
	\loigiai
	{
		Do miền nghiệm không chứa biên nên ta loại đáp án $\heva{& x-2y\le 0 \\& x+3y\ge -2}$ và $\heva{& x-2y\le 0 \\& x+3y\le -2.}$\\
		Chọn điểm $M(0;1)$ thử vào các hệ bất phương trình.\\
		Xét đáp án $\heva{& x-2y<0 \\& x+3y>-2}$, ta có $\heva{& 0-2\cdot 1>0 \\& 0+3\cdot 1>-2 }$ (luôn đúng) và miền nghiệm không chứa biên.
	}
\end{ex}


\begin{ex}%[0D3N1-2]
	Miền nghiệm của hệ bất phương trình  $\heva{& x>0 \\& x-y\le2\\&x+y\le1}$ chứa điểm nào sau đây?
	\choice
	{\True $A\left(\dfrac{1}{2};-1\right)$}
	{$B\left(1;2\right)$}
	{$C\left(0;2\right)$}
	{$D\left(3;-2\right)$}
	\loigiai{
		Thay tọa độ điểm $A\left(\dfrac{1}{2};-1\right)$ vào hệ phương trình  ta thấy thỏa mãn.\\
		Suy ra điểm  $A\left(\dfrac{1}{2};-1\right)$ thuộc miền nghiệm của hệ bất phương trình đã cho.	
	}
\end{ex}

\begin{ex}%[0D3N1-2]
	Cho hệ bất phương trình $\heva{& 2x-5y-1>0 \\& 2x+y+5>0\\&x+y+1<0}$. Trong các điểm sau, điểm nào thuộc miền nghiệm của hệ bất phương trình?
	\choice
	{$O\left(0;0\right)$}
	{$M\left(1;0\right)$}
	{\True $N\left(0;-2\right)$}
	{$P\left(0;2\right)$}
	\loigiai{
		Ta thay lần lượt tọa độ các điểm vào hệ bất phương trình.\\
		Ta thấy điểm $N\left(0;-2\right)$ thỏa mãn hệ.
	}
\end{ex}	

\begin{ex}%[0D2H2-2]
	Miền nghiệm của hệ bất phương trình $\heva{& x+y-1>0\\ & y\geq 2\\ &-x+2y>3}$ là phần \textbf{không} tô đậm của hình vẽ nào trong các hình vẽ sau?
	\choice
	{\begin{tikzpicture}[line join=round, line cap=round, >=stealth,font=\footnotesize, scale=0.6]
			% Vẽ hệ trục toạ độ
			\draw (0,0) node [below left] {$O$};
			\draw [->] (-3.5,0)--(3.5,0) node[below]{$x$};
			\draw[->] (0,-3.5)--(0,3.5) node[left]{$y$};
			\clip(-3.5,-3.5) rectangle (3.5,3.5);
			% Vẽ các số trên hai trục
			\foreach \x in {-3,1}
			\draw[shift={(\x,0)},color=black] (0pt,2pt) -- (0pt,-2pt) node[below] {\footnotesize $\x$};
			\foreach \y in {1,2}
			\draw[shift={(0,\y)},color=black] (2pt,0pt) -- (-2pt,0pt)node[left] {\footnotesize $\y$};
			%Vẽ đồ thị từng hàm số
			\draw[samples=100,smooth,domain=-3.5:3.5,black] plot(\x,{2});
			\draw[samples=100,smooth,domain=-3.5:3.5,black] plot(\x,{1-(\x)});
			\draw[samples=100,smooth,domain=-3.5:3.5,black] plot(\x,{0.5*(\x)+1.5});
			% Vẽ miền nghiệm
			\fill[pattern=north east lines,opacity=0.6] plot[domain=-3.5:3.5](\x,{1-(\x)})--plot[domain=3.5:-3.5](\x,{-3.5})--cycle;
			\fill[pattern=north east lines,opacity=0.6] plot[domain=-3.5:3.5](\x,{2})--plot[domain=3.5:-3.5](\x,{-3.5})--cycle;
			\fill[pattern=north east lines,opacity=0.6] plot[domain=-3.5:3.5](\x,{0.5*(\x)+1.5})--plot[domain=3.5:-3.5](\x,{3.5})--cycle;
	\end{tikzpicture}}
	{\True\begin{tikzpicture}[line join=round, line cap=round, >=stealth,font=\footnotesize, scale=0.6]
			% Vẽ hệ trục toạ độ
			\draw (0,0) node [below left] {$O$};
			\draw [->] (-3.5,0)--(3.5,0) node[below]{$x$};
			\draw[->] (0,-3.5)--(0,3.5) node[left]{$y$};
			\clip(-3.5,-3.5) rectangle (3.5,3.5);
			% Vẽ các số trên hai trục
			\foreach \x in {-3,1}
			\draw[shift={(\x,0)},color=black] (0pt,2pt) -- (0pt,-2pt) node[below] {\footnotesize $\x$};
			\foreach \y in {1,2}
			\draw[shift={(0,\y)},color=black] (2pt,0pt) -- (-2pt,0pt)node[left] {\footnotesize $\y$};
			%Vẽ đồ thị từng hàm số
			\draw[samples=100,smooth,domain=-3.5:3.5,black] plot(\x,{2});
			\draw[samples=100,smooth,domain=-3.5:3.5,black] plot(\x,{1-(\x)});
			\draw[samples=100,smooth,domain=-3.5:3.5,black] plot(\x,{0.5*(\x)+1.5});
			% Vẽ miền nghiệm
			\fill[pattern=north east lines,opacity=0.6] plot[domain=-3.5:3.5](\x,{1-(\x)})--plot[domain=3.5:-3.5](\x,{-3.5})--cycle;
			\fill[pattern=north east lines,opacity=0.6] plot[domain=-3.5:3.5](\x,{2})--plot[domain=3.5:-3.5](\x,{-3.5})--cycle;
			\fill[pattern=north east lines,opacity=0.6] plot[domain=-3.5:3.5](\x,{0.5*(\x)+1.5})--plot[domain=3.5:-3.5](\x,{-3.5})--cycle;
	\end{tikzpicture}}
	{\begin{tikzpicture}[line join=round, line cap=round, >=stealth,font=\footnotesize, scale=0.6]
			% Vẽ hệ trục toạ độ
			\draw (0,0) node [below left] {$O$};
			\draw [->] (-3.5,0)--(3.5,0) node[below]{$x$};
			\draw[->] (0,-3.5)--(0,3.5) node[left]{$y$};
			\clip(-3.5,-3.5) rectangle (3.5,3.5);
			% Vẽ các số trên hai trục
			\foreach \x in {-3,1}
			\draw[shift={(\x,0)},color=black] (0pt,2pt) -- (0pt,-2pt) node[below] {\footnotesize $\x$};
			\foreach \y in {1,2}
			\draw[shift={(0,\y)},color=black] (2pt,0pt) -- (-2pt,0pt)node[left] {\footnotesize $\y$};
			%Vẽ đồ thị từng hàm số
			\draw[samples=100,smooth,domain=-3.5:3.5,black] plot(\x,{2});
			\draw[samples=100,smooth,domain=-3.5:3.5,black] plot(\x,{1-(\x)});
			\draw[samples=100,smooth,domain=-3.5:3.5,black] plot(\x,{0.5*(\x)+1.5});
			% Vẽ miền nghiệm
			\fill[pattern=north east lines,opacity=0.6] plot[domain=-3.5:3.5](\x,{1-(\x)})--plot[domain=3.5:-3.5](\x,{3.5})--cycle;
			\fill[pattern=north east lines,opacity=0.6] plot[domain=-3.5:3.5](\x,{2})--plot[domain=3.5:-3.5](\x,{-3.5})--cycle;
			\fill[pattern=north east lines,opacity=0.6] plot[domain=-3.5:3.5](\x,{0.5*(\x)+1.5})--plot[domain=3.5:-3.5](\x,{-3.5})--cycle;
	\end{tikzpicture}}
	{\begin{tikzpicture}[line join=round, line cap=round, >=stealth,font=\footnotesize, scale=0.6]
			% Vẽ hệ trục toạ độ
			\draw (0,0) node [below left] {$O$};
			\draw [->] (-3.5,0)--(3.5,0) node[below]{$x$};
			\draw[->] (0,-3.5)--(0,3.5) node[left]{$y$};
			\clip(-3.5,-3.5) rectangle (3.5,3.5);
			% Vẽ các số trên hai trục
			\foreach \x in {-3,1}
			\draw[shift={(\x,0)},color=black] (0pt,2pt) -- (0pt,-2pt) node[below] {\footnotesize $\x$};
			\foreach \y in {1,2}
			\draw[shift={(0,\y)},color=black] (2pt,0pt) -- (-2pt,0pt)node[left] {\footnotesize $\y$};
			%Vẽ đồ thị từng hàm số
			\draw[samples=100,smooth,domain=-3.5:3.5,black] plot(\x,{2});
			\draw[samples=100,smooth,domain=-3.5:3.5,black] plot(\x,{1-(\x)});
			\draw[samples=100,smooth,domain=-3.5:3.5,black] plot(\x,{0.5*(\x)+1.5});
			% Vẽ miền nghiệm
			\fill[pattern=north east lines,opacity=0.6] plot[domain=-3.5:3.5](\x,{1-(\x)})--plot[domain=3.5:-3.5](\x,{3.5})--cycle;
			\fill[pattern=north east lines,opacity=0.6] plot[domain=-3.5:3.5](\x,{2})--plot[domain=3.5:-3.5](\x,{3.5})--cycle;
			\fill[pattern=north east lines,opacity=0.6] plot[domain=-3.5:3.5](\x,{0.5*(\x)+1.5})--plot[domain=3.5:-3.5](\x,{-3.5})--cycle;
	\end{tikzpicture}}
	\loigiai{
		Chọn điểm $(0;3)$ là điểm nằm trong miền nghiệm hệ bất phương trình. Do đó ta chọn hình 
		\begin{center}
			\begin{tikzpicture}[line join=round, line cap=round, >=stealth,font=\footnotesize, scale=0.6]
				% Vẽ hệ trục toạ độ
				\draw (0,0) node [below left] {$O$};
				\draw [->] (-3.5,0)--(3.5,0) node[below]{$x$};
				\draw[->] (0,-3.5)--(0,3.5) node[left]{$y$};
				\clip(-3.5,-3.5) rectangle (3.5,3.5);
				% Vẽ các số trên hai trục
				\foreach \x in {-3,1}
				\draw[shift={(\x,0)},color=black] (0pt,2pt) -- (0pt,-2pt) node[below] {\footnotesize $\x$};
				\foreach \y in {1,2}
				\draw[shift={(0,\y)},color=black] (2pt,0pt) -- (-2pt,0pt)node[left] {\footnotesize $\y$};
				%Vẽ đồ thị từng hàm số
				\draw[samples=100,smooth,domain=-3.5:3.5,black] plot(\x,{2});
				\draw[samples=100,smooth,domain=-3.5:3.5,black] plot(\x,{1-(\x)});
				\draw[samples=100,smooth,domain=-3.5:3.5,black] plot(\x,{0.5*(\x)+1.5});
				% Vẽ miền nghiệm
				\fill[pattern=north east lines,opacity=0.6] plot[domain=-3.5:3.5](\x,{1-(\x)})--plot[domain=3.5:-3.5](\x,{-3.5})--cycle;
				\fill[pattern=north east lines,opacity=0.6] plot[domain=-3.5:3.5](\x,{2})--plot[domain=3.5:-3.5](\x,{-3.5})--cycle;
				\fill[pattern=north east lines,opacity=0.6] plot[domain=-3.5:3.5](\x,{0.5*(\x)+1.5})--plot[domain=3.5:-3.5](\x,{-3.5})--cycle;
			\end{tikzpicture}
		\end{center}
	}	
\end{ex}


\begin{ex}%[0D2B2-2]
	Miền tam giác $A B C$ kể cả ba cạnh sau đây là miền nghiệm của hệ bất phương trình nào trong bốn hệ bất phương trình dưới đây?
	\begin{center}
		\begin{tikzpicture}[scale=0.8,thick,>=stealth']
			%\fill[thin,color=gray!50] (-1,3.5) -- (4,3.5) -- (4,-1.2) -- (-1,2.8) -- cycle;
			%\fill[thin,color=gray!50] (-0.4,-3) -- (91/41,10/41) -- (4,-1.2) -- (4,-3) -- cycle;
			%\fill[thin,color=gray!50] (-1,2.8) -- (0,2) -- (0,-3) -- (-0.4,-3) --(-1,-3) -- cycle;
			\fill[thin,color=gray!50] (0,2) -- (0,-2.5) -- (91/41,10/41)  -- cycle;
			\draw[->] (-1,0) -- (4,0)node[above right]{$x$};
			\foreach \x in {1,2,3}
			\draw[shift={(\x,0)},color=black] (0pt,2pt) -- (0pt,-2pt) node[below] {\footnotesize $\x$};
			\draw[->,color=black] (0,-3) -- (0,3.5)node[above right]{$y$};
			\node[below right] at (0,0){$O$};
			\clip(-1,-3) rectangle (4,3.5);
			\draw[line width=1.2pt,smooth,samples=100,domain=-1:4] plot(\x,{5/4*(\x)-10/4});
			\draw[line width=1.2pt,smooth,samples=100,domain=-1:4] plot(\x,{10/5-4/5*(\x)});
			\foreach \y in {1,2,3}
			\draw[shift={(0,\y)},color=black] (2pt,0pt) -- (-2pt,0pt) node[left]{\footnotesize $\y$};
			\path 
			(0,2) coordinate (A)
			(90/41,10/41) coordinate (B)
			(0,-2.5) coordinate (C)
			;
			\foreach \x/\g in {A/60,B/90,C/-10}
			\draw[fill=black] (\x) circle (.036)+(\g:.35)
			node{$\x$};	
			\draw[shift={(2.5,0)},color=black] (0pt,2pt) -- (0pt,-2pt) node[below] {\footnotesize $\frac{5}{2}$};
		\end{tikzpicture}
	\end{center}
	\choice
	{$\heva{
			& y \geq 0 \\
			& 5 x-4 y \geq 10 \\
			& 5 x+4 y \leq 10
		}$}
	{$\heva{
			& x>0 \\
			& 5 x-4 y \leq 10 \\
			& 4 x+5 y \leq 10
		}$} 
	{$\heva{
			& x \geq 0 \\
			& 4 x-5 y \leq 10 \\
			& 5 x+4 y \leq 10
		}$}
	{\True $\heva{
			& x \geq 0 \\
			& 5 x-4 y \leq 10 \\
			& 4 x+5 y \leq 10
		}$}
	\loigiai{Cạnh $A C$ có phương trình $x=0$ và cạnh $A C$ nằm trong miền nghiệm nên $x \geq 0$ là một bất phương trình của hệ.\\[0.15cm]
		Cạnh $A B$ qua hai điểm $\left(\dfrac{5}{2} ; 0\right)$ và $(0 ; 2)$ nên có phương trình $\dfrac{x}{\dfrac{5}{2}}+\dfrac{y}{2}=1 \Leftrightarrow 4 x+5 y=10$.\\[0.15cm]
		Vậy hệ bất phương trình cần tìm là $\heva{
			& x \geq 0 \\
			& 5 x-4 y \leq 10 \\
			& 4 x+5 y \leq 10.
		}$
	}
\end{ex}


\begin{ex}%[0D3N1-2]
	Cho hệ bất phương trình $\heva{&-2x+y\le 2 \\ & -x+2y\ge 4 \\ & x+y\ge 5}$   có miền nghiệm là miền được tô màu (miền tam giác $ABC$  có tọa độ các đỉnh là $A\left(2;0\right)$, $B\left(2;3\right)$, $C\left(1;4\right)$, bao gồm cả các cạnh) như hình vẽ: 
	\begin{center}
		\begin{tikzpicture}[line join=round, line cap=round, >=stealth,font=\footnotesize, scale=0.6]
			% Vẽ hệ trục toạ độ
			\draw (0,0) node [below left] {$O$};
			\draw (0,2) node [below right] {$A$};
			\draw (1,4) node [above] {$B$};
			\draw (2,3) node [above] {$C$};
			\draw [->] (-2,0)--(4.5,0) node[below]{$x$};
			\draw[->] (0,-1)--(0,6) node[left]{$y$};
			\clip(-2,-1) rectangle (4.5,6);
			% Vẽ các số trên hai trục
			\foreach \x in {1,2}
			\draw[shift={(\x,0)},color=black] (0pt,2pt) -- (0pt,-2pt) node[below] {\footnotesize $\x$};
			\foreach \y in {2,3,4,5}
			\draw[shift={(0,\y)},color=black] (2pt,0pt) -- (-2pt,0pt)node[left] {\footnotesize $\y$};
			%Vẽ đồ thị từng hàm số
			\draw[samples=100,smooth,domain=-6.5:4.5,black] plot(\x,{2*(\x)+2});
			\draw[samples=100,smooth,domain=-6.5:4.5,black] plot(\x,{0.5*(\x)+2});
			\draw[samples=100,smooth,domain=-6.5:4.5,black] plot(\x,{-(\x)+5});
			% Vẽ miền nghiệm
			\fill[pattern=north west lines,opacity=0.6] (0,2)--(2,3)--(1,4)--cycle;
			\draw [dashed] (1,0)--(1,4)--(0,4) (2,0)--(2,3)--(0,3) ;
		\end{tikzpicture}
	\end{center}
	Giá trị nhỏ nhất $F_{\min}$  của biểu thức $F\left(x; y\right)$  trên miền xác định bởi hệ trên là 
	\choice
	{\True $F_{\min}=1$}
	{$F_{\min}=2$}
	{$F_{\min}=3$}
	{$F_{\min}=4$ }
	\loigiai{
		Tại $A\left(2;0\right)$ ta có  $F\left(2;0\right)=2$.\\
		Tại $B\left(2;3\right)$  ta có $F\left(2;3\right)=1$.\\
		Tại $C\left(1;4\right)$  ta có  $F\left(1;4\right)=3$.\\
		Vậy  $F_{\min}=1$.		
	}
\end{ex}

\begin{ex}%[0D3N1-2]
	Giá trị nhỏ nhất $F_{\min}$   của biểu thức  $F\left(x; y\right)$  trên miền xác định bởi hệ  $\heva{&0 \leq x \leq 10\\&0 \leq y \leq 9\\&2x+y\geq 14\\&2x+5y\geq 30} $ là
	\choice
	{$F_{\min}=23$}
	{$F_{\min}=26$}
	{\True $F_{\min}=32$}
	{$F_{\min}=67$}
	\loigiai{
		Trong mặt phẳng tọa độ  vẽ các đường thẳng $d_1:2x+y-14=0$, $d_1:2x+5y-30=0$, $\Delta:y=9$, $\Delta'\colon x=10$.\\		
		Khi đó miền nghiệm của hệ bất phương trình là phần mặt phẳng  tô màu như hình vẽ.
		\begin{center}
			\begin{tikzpicture}[scale=0.6,thick,>=stealth']
				\draw[->] (-1,0) -- (15.3,0)node[above]{$x$};
				\foreach \x in {1,2,3,4,5,6,7,8,9,10,11,12,13,14,15}
				\draw[shift={(\x,0)},color=black] (0pt,2pt) -- (0pt,-2pt) node[below] {\footnotesize $\x$};
				\draw[->,color=black] (0,-1) -- (0,14.3)node[right]{$y$};
				\foreach \y in {1,2,3,4,5,6,7,8,9,10,11,12,13,14}
				\draw[shift={(0,\y)},color=black] (2pt,0pt) -- (-2pt,0pt) node[left] {\footnotesize $\y$};
				\node[below left] at (0,0) {$O$};
				\clip(-1,-1) rectangle (15.3,14.3);
				%\fill[pattern=north east lines] (-0.15,14.3) -- (-1,14.3) -- (-1,-1) -- (7.5,-1)-- cycle;
				\draw[line width=1.2pt,smooth,samples=100,domain=-1:9] plot(\x,{14-2*(\x)});
				%\fill[pattern=north east lines] (-1,6.4)--(-1,-1)--(15.3,-1)--(15.3,-0.12)-- cycle;
				\draw[line width=1.2pt,smooth,samples=100,domain=-1:15.3] plot(\x,{6-0.4*(\x)});
				%\fill[pattern=north east lines] (-1,9)--(-1,14.3) --(15.3,14.3)--(15.3,9)-- cycle;
				\draw[line width=1.2pt,smooth,samples=100,domain=-1:15.3] plot(\x,{9+0*(\x)});
				%\fill[pattern=north east lines] (10,14.3)--(15.3,14.3) --(15.3,-1)--(10,-1)-- cycle;
				\draw[line width=1.2pt,smooth,samples=100](10,-1)--(10,15.3);
				%\fill[pattern=north east lines] (0,4)--(-3,4) --(-3,-3)--(0,-3)-- cycle;
				\draw[line width=1.2pt,smooth,samples=100,domain=-3:9] plot(\x,{0*(\x)});
				\fill[pattern=north east lines] (5,4)--(10,2) --(10,9)--(2.5,9)-- cycle;
				\node[below left] at (5.2,4.2) {$A$};
				\node[right] at (10,3) {$B$};
				\node[right] at (10,8.5) {$C$};
				\node[left] at (2.75,8.5) {$D$};	
			\end{tikzpicture}
		\end{center}
		Xét các đỉnh của miền khép kín tạo bởi hệ là  $A\left(5; 4\right)$, $B\left(\dfrac{5}{2}; 9\right)$, $C\left(10; 9\right)$, $D\left(10; 2\right)$.\\
		Ta có   $F\left(5;4\right)=32$.\\
		$F\left(\dfrac{5}{2}; 9\right)=37$ .\\
		$F\left(10;9\right)=67$.\\
		$F\left(10;2\right)=46$.\\
		Vậy  $F_{\min}=32$.	
	}
\end{ex}	

\begin{ex}%[0D3N1-2]
	Giá trị lớn nhất $F_{\max}$  của biểu thức $F\left(x; y\right)=x+2y$  trên miền xác định bởi hệ  $\heva{&0 \leq y \leq 4\\&x\geq0\\&x-y-1\leq 0\\&x+2y-10\leq 0} $ là
	\choice
	{$F_{\max}=6$ }
	{$F_{\max}=8$}
	{\True $F_{\max}=10$}
	{$F_{\max}=12$}
	\loigiai{
		Trong mặt phẳng tọa độ  vẽ các đường thẳng $d_1:x-y-1=0$, $d_2:x+2y-10=0$, $\Delta:y=4$ \\
		Khi đó miền nghiệm của hệ bất phương trình là phần mặt phẳng  tô màu như hình vẽ.
		\begin{center}
			\begin{tikzpicture}[scale=0.6,thick,>=stealth']
				\draw[->] (-3,0) -- (12,0)node[above]{$x$};
				\foreach \x in {1,2,3,4,5,6,7,8,9,10,11}
				\draw[shift={(\x,0)},color=black] (0pt,2pt) -- (0pt,-2pt) node[below] {\footnotesize $\x$};
				\draw[->,color=black] (0,-2) -- (0,8)node[right]{$y$};
				\foreach \y in {1,2,3,4,5,6,7}
				\draw[shift={(0,\y)},color=black] (2pt,0pt) -- (-2pt,0pt) node[left] {\footnotesize $\y$};
				\node[below left] at (0,0) {$O$};
				\clip(-3,-2) rectangle (12,8);
				%\fill[pattern=north east lines] (-0.15,14.3) -- (-1,14.3) -- (-1,-1) -- (7.5,-1)-- cycle;
				\draw[line width=1.2pt,smooth,samples=100,domain=-1:9] plot(\x,{(\x)-1});
				%\fill[pattern=north east lines] (-1,6.4)--(-1,-1)--(15.3,-1)--(15.3,-0.12)-- cycle;
				\draw[line width=1.2pt,smooth,samples=100,domain=-1:15.3] plot(\x,{-0.5*(\x)+5});
				%\fill[pattern=north east lines] (-1,9)--(-1,14.3) --(15.3,14.3)--(15.3,9)-- cycle;
				\draw[line width=1.2pt,smooth,samples=100,domain=-1:15.3] plot(\x,{4});
				
				\fill[pattern=north east lines] (0,0)--(0,4) --(2,4)--(4,3)--(1,0)-- cycle;
				\node[above] at (1,0) {$A$};
				\node[right] at (4,3) {$B$};
				\node[above] at (2,4) {$C$};
				\node[below right] at (0,4) {$D$};	
			\end{tikzpicture}
		\end{center}
		Xét các đỉnh của miền khép kín tạo bởi hệ là $O\left(0;0\right)$,  $A\left(1;0\right)$,  $B\left(4;3\right)$,  $C\left(2;4\right)$,  $D\left(0;4\right)$\\
		Ta có  $\heva{&F\left(0; 0\right)=0\\ &F\left(1; 0\right)=1\\&F\left(4; 3\right)=10\\&F\left(2; 4\right)=10\\&F\left(0; 4\right)=4} $ $\Rightarrow F_{\max}=10$.		
	}
\end{ex}
\Closesolutionfile{ans}
\indapan{6}{ans/ans-0-De2}
\cauds
\Opensolutionfile{ans}[ans/ans-0-De2-DS]
\begin{ex}%[0D3H1-5]
	Cho hệ bất phương trình $\heva{&  x+ y -2\leq 0 \\& 2 x-3 y +2>0}$. Các mệnh đề sau đúng hay sai?
	\choiceTF
	{\True Hệ trên là một hệ bất phương trình bậc nhất hai ẩn}
	{$\left(-1; 3\right)$ là một nghiệm của hệ bất phương trình trên}
	{\True $\left(1; 1\right)$ là một nghiệm của hệ bất phương trình trên}
	{$\left(2; -2\right)$  không là một nghiệm của hệ bất phương trình trên}
	\loigiai{
		\begin{itemchoice}
			\itemch Đúng vì hệ trên là một hệ bất phương trình bậc nhất hai ẩn. 
			\itemch Sai vì thay $x=-1; y=3$  vào hệ bất phương trình ta được $\heva{&0\ge0\\&-9>0}$ (Vô lý).\\ Vậy  $\left(-1; 3\right)$  không phải là một nghiệm của hệ bất phương trìnhai.
			\itemch Đúng vì thay $x=1; y=1$   vào hệ bất phương trình ta được $\heva{&0\le0\\&1>0}$ (luôn đúng).\\ Vậy $\left(1; 1\right)$  là một nghiệm của hệ bất phương trình.
			\itemch Sai vì thay $x=2; y=-2$  vào hệ bất phương trình ta được $\heva{&-2\le0\\&12>0}$ (luôn đúng).\\ Vậy $\left(2; -2\right)$  là một nghiệm của hệ bất phương trình.
		\end{itemchoice}
	}
\end{ex}

\begin{ex}%[0D3H1-5]
	Cho hệ bất phương trình $\heva{&x+2y-100 \leq 0\\&2x+y-80\leq0\\&x\geq 0\\&y\geq 0}$. Các mệnh đề sau đúng hay sai?	
	\choiceTF
	{Hệ trên không phải là một hệ bất phương trình bậc nhất hai ẩn}
	{$\left(-1; 2\right)$   là một nghiệm của hệ bất phương trình trên}
	{\True $\left(-1; 5\right)$ không là một nghiệm của hệ bất phương trình trên}
	{\True Miền nghiệm của hệ bất phương trình trên là một miền tứ giác}
	\loigiai{
		\begin{itemchoice}
			\itemch Sai. Vì hệ trên là một hệ bất phương trình bậc nhất hai ẩn.
			\itemch Sai.  Vì thay $x=-1; y=2$  vào hệ bất phương trình ta được $\heva{-97 &\leq 0\\-80&\leq0\\-1&\geq 0\\2&\geq 0} $  (Vô lý).\\ Vậy $\left(-1; 2\right)$  không phải là một nghiệm của hệ bất phương trình.
			\itemch Đúng vì thay  $x=-1; y=5$ vào hệ bất phương trình ta được $\heva{&-91 \leq 0\\&-77\leq0\\&-1\geq 0\\&5\geq 0} $  (Vô lý).\\ Vậy $\left(-1; 5\right)$  không là một nghiệm của hệ bất phương trình.
			\itemch Đúng vì vẽ đường thẳng $d_1:x+2y-100=0$. Vì $0+2\cdot0-100\leq0$  nên toạ độ điểm $O\left(0; 0\right)$   thoả mãn bất phương trình $:x+2y-100\leq0$. Do đó miền nghiệm $D_1$  của bất phương trình $x+2y-100\leq0$ là nửa mặt phẳng bờ $d_1:x+2y-100=0$  chứa gốc toạ độ $O$.\\ 
			Tương tự miền nghiệm $D_2$   của bất phương trình $2x+y-80\leq0$ là nửa mặt phẳng bờ $d_2:2x+y-80=0$  chứa gốc toạ độ $O$.\\  
			Tương tự miền nghiệm  $D_3$  của bất phương trình $x\geq0$ là nửa mặt phẳng bờ  $Oy$ chứa điểm  $\left(1; 0\right)$.\\
			Tương tự miền nghiệm $D_3$  của bất phương trình $y\geq0$ là nửa mặt phẳng bờ $Ox$  chứa điểm  $\left(0; 1\right)$.\\
			Khi đó miền nghiệm của hệ bất phương trình là miền tứ giác $OABC$. 
			\begin{center}
				\begin{tikzpicture}[scale=0.6,thick,>=stealth']
					\draw[->] (-3,0) -- (12,0)node[above]{$x$};
					%\foreach \x in {1,2,3,4,5,6,7,8,9,10,11,11,12}\draw[shift={(\x,0)},color=black] (0pt,2pt) -- (0pt,-2pt) node[below] {\footnotesize $\x$};
					\draw[->,color=black] (0,-2) -- (0,9)node[right]{$y$};
					%\foreach \y in {1,2,3,4,5,6,7,8,9}\draw[shift={(0,\y)},color=black] (2pt,0pt) -- (-2pt,0pt) node[left] {\footnotesize $\y$};
					\node[below left] at (0,0) {$O$};
					\clip(-3,-2) rectangle (12,9);
					%\fill[pattern=north east lines] (-0.15,14.3) -- (-1,14.3) -- (-1,-1) -- (7.5,-1)-- cycle;
					\draw[line width=1.2pt,smooth,samples=100,domain=-1:9] plot(\x,{-2*(\x)+8});
					%\fill[pattern=north east lines] (-1,6.4)--(-1,-1)--(15.3,-1)--(15.3,-0.12)-- cycle;
					\draw[line width=1.2pt,smooth,samples=100,domain=-1:15.3] plot(\x,{-0.5*(\x)+5});
					%\fill[pattern=north east lines] (-1,9)--(-1,14.3) --(15.3,14.3)--(15.3,9)-- cycle;
					
					\fill[pattern=north east lines] (0,0)--(0,5) --(2,4)--(4,0)-- cycle;
					\node[above] at (4,0) {$A$};
					\node[below] at (4,0) {$40$};\node[below] at (10,0) {$100$};
					\node[right] at (0,5) {$B$}; \node[left] at (0,5) {$50$};
					\node[left] at (0,8) {$80$};
					\node[right] at (2,4) {$C$};\node[left] at (0,4) {$40$};
					\node[below] at (2,0) {$20$};\draw[dashed] (0,4)--(2,4)--(2,0);
				\end{tikzpicture}
			\end{center}
		\end{itemchoice}
	}
\end{ex}

\begin{ex}%[0D3H1-5]
	Cho hệ bất phương trình $\heva{	& 2 x+3 y -1>0 \\& 5 x+ y-4 <0}$?	
	\choiceTF
	{\True Hệ trên là một hệ bất phương trình bậc nhất hai ẩn}
	{$\left(1; 2\right)$ là một nghiệm của hệ bất phương trình trên}
	{Điểm  $A\left(-2; 1\right)$ thuộc miền  nghiệm của hệ bất phương trình trên}
	{Hệ bất phương trình trên vô nghiệm.}
	\loigiai{
		\begin{itemchoice}
			\itemch Đúng. Vì hệ trên là một hệ bất phương trình bậc nhất hai ẩn.
			\itemch Sai. Vì thay $x=1; y=2$  vào hệ bất phương trình ta được $\heva{
				& 7>0 \\
				& 7 <0
			}$ (Vô lý).\\ Vậy $\left(1; 2\right)$  không phải là một nghiệm của hệ bất phương trình.
			\itemch Sai. Vì thay   vào hệ bất phương trình ta được  (Vô lý). Vậy  điểm   không thuộc miền nghiệm của hệ bất phương trình.
			\itemch Sai. Vì hệ bất phương trình trên có nghiệm $\left(-1; 1\right)$.
		\end{itemchoice}
	}
\end{ex}
\begin{ex}%[0D3N1-5]
	 Xét tính đúng, sai của các mệnh đề sau
	 \choiceTF
{$\heva{&2 x+y \leq 3 \\& 3 x-2 y \geq 1}$ là hệ bất phương trình bậc nhất hai ẩn}
{$\heva{&x+2 \leq 0 \\& y-3 \geq 1}$ là hệ bất phương trình bậc nhất hai ẩn}
{$\heva{&x+2 y=2 \\& 2 x-3 y=-1}$ là hệ bất phương trình bậc nhất hai ẩn}
{$\heva{&2 x+5 y \leq 30 \\&4 x-3 y \geq 5 \\ x \geq 0 \\ y \geq 1}$ là hệ bất phương trình bậc nhất hai ẩn}
\begin{itemchoice}
\itemch Đúng. $\heva{&2 x+y \leq 3 \\& 3 x-2 y \geq 1}$ là hệ bất phương trình bậc nhất hai ẩn.
\itemch Đúng. $\heva{&x+2 \leq 0 \\& y-3 \geq 1}$ là hệ bất phương trình bậc nhất hai ẩn.
\itemch Sai. Vì $\heva{&x+2 y=2 \\& 2 x-3 y=-1}$ là hệ  phương trình bậc nhất hai ẩn.
\itemch Đúng. $\heva{&2 x+5 y \leq 30 \\&4 x-3 y \geq 5 \\ x \geq 0 \\ y \geq 1}$ là hệ bất phương trình bậc nhất hai ẩn.
\end{itemchoice}
	
\end{ex}	
\begin{ex}%[0D3H1-5]
	Cho hệ bất phương trình $\heva{ & 2x-\frac{3}{2}y\ge 1 \qquad (1)\\ & 4x-3y\ge 2 \qquad (2)}$ có tập nghiệm $S$. Mệnh đề nào sau đây là đúng?	
	\choiceTF
	{\True Hệ bất phương trình trên là một hệ bất phương trình bậc nhất hai ẩn}
	{$\left(-\frac{1}{4}; -1\right)\notin S$  }
	{Điểm $A\left(2; 3\right)$  thuộc miền nghiệm của hệ bất phương trình đã cho}
	{\True $S=\{\left(x; y\right)|4x-3y=2\}$}
	\loigiai{
		\begin{itemchoice}
			\itemch Đúng. Vì hệ bất phương trình trên là một hệ bất phương trình bậc nhất hai ẩn nên mệnh đề đúng;
			\itemch Sai. Vì thay  $x=-\frac{1}{4}; y=-1$ vào hệ bất phương trình ta được $\heva{ & 1\ge 1 \\ & 2\ge 2}$ (đúng).\\ Vậy $\left(-\frac{1}{4}; -1\right)$  là    
			một nghiệm của hệ bất phương trình nên mệnh đề sai.
			
			\itemch Sai. Vì thay $x=2; y=3$  vào hệ bất phương trình ta được $\heva{ & -\frac{1}{2}\ge 1 \\ & -1\ge 2}$ (Vô lý).\\ Vậy $A\left(2; 3\right)$  không thuộc miền nghiệm của hệ bất phương trình nên mệnh đề sai.
			\itemch Đúng. Vì ta vẽ hai đường thẳng: $d_1:2x-\dfrac{3}{2}y=1$, $d_2:4x-3y=2$
			\begin{center}
				\begin{tikzpicture}[scale=0.6,thick,>=stealth']
					\fill[pattern=north east lines] (-3,-3)--(-3,6) --(6,6)--(6,-3)-- cycle;
					\draw[->] (-3,0) -- (6,0)node[above]{$x$};
					\node[above left] at (0,0) {$O$};
					\clip(-3,-3) rectangle (6,6);
					\draw[->,color=black] (0,-3) -- (0,6)node[right]{$y$};
					\draw[line width=1.2pt,smooth,samples=100,domain=-3:6] plot(\x,{(4/3)*(\x)-2/3});
					\node[above] at (.5,0) {$\frac{1}{2}$};
					\node[right] at (0,-2/3) {$-\frac{2}{3}$};
				\end{tikzpicture}
			\end{center}
			Thử trực tiếp ta thấy $\left(0; 0\right)$  là nghiệm của bất phương trình $(2)$ nhưng không phải là nghiệm         
			của bất phương trình $(1)$. Sau khi gạch bỏ các miền không thích hợp, tập hợp nghiệm của hệ         
			bất phương trình chính là các điểm thuộc đường thẳng $d:4x-3y=2$  nên mệnh đề đúng.			
		\end{itemchoice}
	}
\end{ex}
\Closesolutionfile{ans}
\indapan{3}{ans/ans-0-De2-DS}

\Opensolutionfile{ans}[ans/ans-0-De2-KQ]
\caukq
\begin{ex}%[0D3N1-3]
	Giá trị nhỏ nhất của biểu thức $F\left(x; y\right)=x-2y$  với điều kiện  $\heva{&0 \leq y \leq 5\\&x \leq 0\\&x+y-2\geq 0\\&x-y-2\leq 0} $   là
	\shortans{$-10$}
	\loigiai{
		Vẽ các đường thẳng $d_1: y=5$, $d_2: x+y-2=0$, $d_3: x-y-2=0$, $Ox: y=0$, $Oy: x=0$.
		\begin{center}
			\begin{tikzpicture}[scale=0.6,thick,>=stealth']
				\draw[->] (-2,0) -- (9,0)node[above]{$x$};
				\foreach \x in {1,2,3,4,5,6,7,8}
				\draw[shift={(\x,0)},color=black] (0pt,2pt) -- (0pt,-2pt) node[below] {\footnotesize $\x$};
				\draw[->,color=black] (0,-3) -- (0,7)node[right]{$y$};
				\foreach \y in {1,2,3,4,5,6,7}
				\draw[shift={(0,\y)},color=black] (2pt,0pt) -- (-2pt,0pt) node[left] {\footnotesize $\y$};
				\node[below left] at (0,0) {$O$};
				\clip(-2,-3) rectangle (9,7);
				%\fill[pattern=north east lines] (-0.15,14.3) -- (-1,14.3) -- (-1,-1) -- (7.5,-1)-- cycle;
				\draw[line width=1.2pt,smooth,samples=100,domain=-1:9] plot(\x,{-(\x)+2});
				%\fill[pattern=north east lines] (-1,6.4)--(-1,-1)--(15.3,-1)--(15.3,-0.12)-- cycle;
				\draw[line width=1.2pt,smooth,samples=100,domain=-1:15.3] plot(\x,{(\x)-2});
				%\fill[pattern=north east lines] (-1,9)--(-1,14.3) --(15.3,14.3)--(15.3,9)-- cycle;
				\draw[line width=1.2pt,smooth,samples=100,domain=-2:8] plot(\x,{5});
				
				\fill[pattern=north east lines] (0,2)--(0,5) --(7,5)--(2,0)-- cycle;
				\node[below right] at (0,5) {$A$};
				\node[right] at (0,2) {$B$};
				\node[above] at (2,0) {$C$};
				\node[below right] at (7,5) {$D$};	
			\end{tikzpicture}
		\end{center}
		Các đường thẳng trên đôi một cắt nhau tại  $A\left(0; 5\right)$, $B\left(0; 2\right)$, $C\left(2; 0\right)$, $D\left(7; 5\right)$.\\ 
		Vì điểm $M_0\left(2; 1\right)$ có toạ độ thoả mãn tất cả các bất phương trình trong hệ nên ta tô đậm các nửa mặt phẳng bờ $d_1$, $d_2$, $d_3$, $Ox$, $Oy$  chứa điểm $M_0\left(2; 1\right)$. Miền  bị tô đậm là đa giác $ABCD$  kể cả các cạnh (hình bên) là miền nghiệm của hệ phương trình đã cho.\\
		Kí hiệu  $F\left(A\right)=F\left(x_A; y_A\right)=x_A-2y_A$.\\
		Ta có
		$F\left(A\right)=-10$, $F\left(B\right)=-4$, $F\left(C\right)=2$, $F\left(D\right)=-3$.\\
		Vậy giá trị nhỏ nhất cần tìm là $F\left(A\right)=-10$.	
	}
\end{ex}

\begin{ex}%[0D3N1-3]
	Biểu thức $L=y-x$, với $x$  và $y$  thỏa mãn hệ bất phương trình $\heva{&2x+3y-6 \leq 0\\&x\geq 0\\&2x-3y-1\leq 0}$, đạt giá trị lớn nhất 
	là $a$  và đạt giá trị nhỏ nhất là $b$. Khi đó $a+12b$ bằng bao nhiêu?
	\shortans{$-9$}
	\loigiai{
		Trước hết, ta vẽ ba đường thẳng $d_1: 2x+3y-6=0$, $d_2:x=0$, $d_3: 2x-3y-1=0$
		\begin{center}
			\begin{tikzpicture}[line join=round, line cap=round, >=stealth,font=\footnotesize, scale=0.6]
				% Vẽ hệ trục toạ độ
				\draw (0,0) node [above left] {$O$};
				\draw [->] (-3.5,0)--(3.5,0) node[below]{$x$};
				\draw[->] (0,-3.5)--(0,3.5) node[left]{$y$};
				\clip(-3.5,-3.5) rectangle (3.5,3.5);
				% Vẽ các số trên hai trục
				
				%Vẽ đồ thị từng hàm số
				\draw[samples=100,smooth,domain=-3.5:3.5,black] plot(\x,{-(2/3)*(\x)+2});
				\draw[samples=100,smooth,domain=-3.5:3.5,black] plot(\x,{(2/3)*(\x)-(1/3)});
				% Vẽ miền nghiệm
				\fill[pattern=north east lines,opacity=0.6] plot[domain=-3.5:3.5](\x,{-(2/3)*(\x)+2})--plot[domain=3.5:-3.5](\x,{3.5})--cycle;
				
				\fill[pattern=north east lines,opacity=0.6] plot[domain=-3.5:3.5](\x,{(2/3)*(\x)-(1/3)})--plot[domain=3.5:-3.5](\x,{-3.5})--cycle;
				\fill[pattern=north east lines,opacity=0.6] (-3.5,-3.5)--(0,-3.5)--(0,3.5)--(-3.5,3.5)--cycle;
				\node[right] at (0,2) {$A$};
				\node[below] at (7/4,5/6) {$B$};
				\node[right] at (0,-1/3) {$C$};
			\end{tikzpicture}
		\end{center}	
		Ta thấy $\left(0;0\right)$  là nghiệm của cả ba bất phương trình. Điều đó có nghĩa gốc tọa độ thuộc cả ba miền nghiệm của cả ba bất phương trình. Sau khi gạch bỏ các miền không thích hợp, miền không bị gạch là miền nghiệm của hệ (kể cả biên).
		Miền nghiệm là hình tam giác $ABC$  (kể cả biên), với  $A\left(0;2\right)$,  $B\left(\dfrac{7}{4};\dfrac{5}{6}\right)$,  $C\left(0;-\dfrac{1}{3}\right)$.\\
		Ta có   $L\left(A\right)=L\left(2;0\right)=2-0=2$,\\ $L\left(B\right)=L\left(\dfrac{7}{4};\dfrac{5}{6} \right)=\dfrac{6}{5}-\dfrac{7}{4}=-\dfrac{11}{12}$,\\ $L\left(C\right)=L\left(0;-\dfrac{1}{3} \right)=-\dfrac{1}{3}-0=-\dfrac{1}{3}$.\\
		Vậy $a=2$, $b=-\dfrac{11}{12}$.\\
		Khi đó $a+12b=-9$.  	
	}
\end{ex}

\begin{ex}%[0D3N1-3]
	Cho các giá trị  $x$, $y$ thỏa mãn điều kiện $\heva{&x-y+2 \geq 0\\&2x-y-1\leq 0\\&3x-y-2\geq 0}$. Tìm giá trị lớn nhất của biểu thức $T=3x+2y$.
	\shortans{$19$}
	\loigiai{
		\begin{center}
			\begin{tikzpicture}[line join=round, line cap=round, >=stealth,font=\footnotesize, scale=0.6,yscale=0.6]
				\draw [pattern=north west lines] (-4,-5)--(-4,10)--(5,10)--(5,-5)--cycle;
				\fill[white] (1,1)--(2,4) --(3,5)-- cycle;
				
				% Vẽ hệ trục toạ độ
				\draw (0,0) node [above left] {$O$};
				\draw [->] (-4,0)--(5,0) node[below]{$x$};
				\draw[->] (0,-5)--(0,10) node[left]{$y$};
				\clip(-4,-5) rectangle (5,10);
				% Vẽ các số trên hai trục
				
				%Vẽ đồ thị từng hàm số
				\draw[samples=100,smooth,domain=-4:5,black] plot(\x,{(\x)+2});
				\draw[samples=100,smooth,domain=-4:5,black] plot(\x,{2*(\x)-1});
				
				
				\draw[samples=100,smooth,domain=-4:5,black] plot(\x,{3*(\x)-2});
				% Vẽ miền nghiệm
				
				
				\node[right] at (1,1) {$A$};
				\node[left] at (2,4) {$B$};
				\node[right] at (3,5) {$C$};
			\end{tikzpicture}
		\end{center}
		Miền nghiệm của hệ đã cho là miền trong tam giác $ABC$ (Kể cả đường biên) trong đó $A\left(1;1\right)$,  $B\left(2;4\right)$,  $C\left(3;5\right)$.
		Giá trị lớn nhất của $T=3x+2y$  đạt được tại các đỉnh của tam giác $ABC$.\\
		Ta có   $T\left(A\right)=T\left(1;1\right)=5$, $T\left(B\right)=T\left(2;4 \right)=14$, $T\left(C\right)=T\left(3;5\right)=19$.\\
		Vậy giá	trị lớn nhất của biểu thức $T=3x+2y=19$.
	}
\end{ex}

\begin{ex}%[0D3N1-3]
	Biết miền nghiệm của hệ bất phương trình $\heva{&3x+y\leq 6\\&x+y\leq 4\\&x,\,y\geq 0}$ là miền tứ giác $OABC$  như hình vẽ dưới đây
	\begin{center}
		\begin{tikzpicture}[scale=0.6,thick,>=stealth']
			\draw[->] (-3,0) -- (7,0)node[above]{$x$};
			\foreach \x in {1,2,3,4,5,6,}
			\draw[shift={(\x,0)},color=black] (0pt,2pt) -- (0pt,-2pt) node[below] {\footnotesize $\x$};
			\draw[->,color=black] (0,-2) -- (0,8)node[right]{$y$};
			\foreach \y in {1,2,3,4,5,6,7}
			\draw[shift={(0,\y)},color=black] (2pt,0pt) -- (-2pt,0pt) node[left] {\footnotesize $\y$};
			\node[below left] at (0,0) {$O$};
			\clip(-3,-2) rectangle (12,8);
			%\fill[pattern=north east lines] (-0.15,14.3) -- (-1,14.3) -- (-1,-1) -- (7.5,-1)-- cycle;
			\draw[line width=1.2pt,smooth,samples=100,domain=-1:9] plot(\x,{-3*(\x)+6});
			%\fill[pattern=north east lines] (-1,6.4)--(-1,-1)--(15.3,-1)--(15.3,-0.12)-- cycle;
			\draw[line width=1.2pt,smooth,samples=100,domain=-1:15.3] plot(\x,{-(\x)+4});
			
			\fill[pattern=north east lines] (0,0)--(0,4) --(1,3)--(2,0)-- cycle;
			\node[above] at (0,4) {$A$};
			\node[right] at (1,3) {$B$};
			\node[above] at (2,0) {$C$};
			
		\end{tikzpicture}
	\end{center}
	Gọi $\left(x_0;y_0\right)$  là một nghiệm của hệ bất phương trình trên sao cho biểu thức $P=2x+y$  đạt giá trị lớn nhất. Khi đó tổng $T=3x_0+7y_0$  bằng bao nhiêu?
	\shortans{$24$}
	\loigiai{
		Ta có: $O\left(0;0\right)$, $A\left(0;4\right)$,  $B\left(1;3\right)$,  $C\left(2;0\right)$. \\
		Biểu thức $P=2x+y$  đạt giá trị lớn nhất tại một trong các đỉnh của tứ giác $OABC$.\\
		Tại  $O\left(0;0\right)$, ta có $P=2x+y=0$ \\
		Tại  $A\left(0;4\right)$, ta có $P=2x+y=4$ \\
		Tại  $B\left(1;3\right)$, ta có  $P=2x+y=5$\\
		Tại  $C\left(2;0\right)$, ta có  $P=2x+y=4$\\
		Ta thấy biểu thức $P=2x+y$  đạt giá trị lớn nhất tại $B\left(1;3\right)$.\\
		Vậy $T=3x_0+7y_0=3\cdot1+7\cdot3=24$.	
	}
\end{ex}

\begin{ex}%[0D3N1-3]
	Một công ty điện tử sản xuất hai kiểu radio trên hai dây chuyền độc lập. Radio kiểu một sản xuất trên dây chuyền một với công suất  $45$ radio/ngày, radio kiểu hai sản xuất trên dây chuyền hai với công suất $80$  radio/ngày. Để sản xuất một chiếc radio kiểu một cần $12$  linh kiện, để sản xuất một chiếc radio kiểu hai cần $9$ linh kiện. Tiền lãi khi bán một chiếc radio kiểu một là $250\,000$  đồng, lãi thu được khi bán một chiếc radio kiểu hai là $180\,000$ đồng. Biết rằng số linh kiện có thể sử dụng tối đa trong một ngày là $900$. Gọi $x_0$, $y_0$  lần lượt là số radio kiểu một và radio kiểu hai sản suất được trong một ngày để tiền lãi thu được là nhiều nhất. Khi đó tổng $T=x_0+2y_0$  bằng
	\shortans{$125$}
	\loigiai{
		Gọi $x$  và  $y$ lần lượt là số radio kiểu một và số radio kiểu hai mà công ty này sản xuất trong một ngày $\left(x,\, y \in \mathbb{N}^*\right)$.\\ 
		Số tiền lãi mà công ty này thu về hàng ngày là $P=250\,000x+180\,000y$  (đồng).\\
		Ta có hệ bất phương trình  $\heva{&12x+9y\leq 900\\&0\leq x\leq 45\\&0\leq y\leq 80.}$\\
		Bài toán trở thành tìm giá trị lớn nhất của biểu thức $P=250\,000x+180\,000y$  trên miền nghiệm của hệ bất phương trình $\heva{&12x+9y\leq 900\\&0\leq x\leq 45\\&0\leq y\leq 80.}$
		\begin{center}
			\begin{tikzpicture}[scale=0.6,thick,>=stealth']
				\draw[->] (-3,0) -- (10,0)node[above]{$x$};
				%\foreach \x in {1,2,3,4,5,6,7,8,9,10,11,11,12}\draw[shift={(\x,0)},color=black] (0pt,2pt) -- (0pt,-2pt) node[below] {\footnotesize $\x$};
				\draw[->,color=black] (0,-2) -- (0,11)node[right]{$y$};
				%\foreach \y in {1,2,3,4,5,6,7,8,9}\draw[shift={(0,\y)},color=black] (2pt,0pt) -- (-2pt,0pt) node[left] {\footnotesize $\y$};
				\node[below left] at (0,0) {$O$};
				\clip(-3,-2) rectangle (10,11);
				%\fill[pattern=north east lines] (-0.15,14.3) -- (-1,14.3) -- (-1,-1) -- (7.5,-1)-- cycle;
				\draw[line width=1.2pt,smooth,samples=100,domain=-1:9] plot(\x,{-(12/9)*(\x)+10});
				%\fill[pattern=north east lines] (-1,6.4)--(-1,-1)--(15.3,-1)--(15.3,-0.12)-- cycle;
				\draw[line width=1.2pt,smooth,samples=100,domain=-1:15.3] plot(\x,{8});
				%\fill[pattern=north east lines] (-1,9)--(-1,14.3) --(15.3,14.3)--(15.3,9)-- cycle;
				
				\fill[pattern=north east lines] (0,0)--(0,8) --(1.5,8)--(4.5,4)-- (4.5,0)-- cycle;
				\node[above right] at (4.5,0) {$A$};
				\node[below right] at (4.5,0) {$45$};\node[below] at (7.5,0) {$75$};
				\node[right] at (4.5,4) {$B$}; \node[below] at (1.5,0) {$15$};
				\node[below left] at (0,8) {$80$};
				\node[above right] at (1.5,8) {$C$};\node[left] at (0,4) {$40$};
				\node[left] at (0,10) {$100$};
				\node[above left] at (0,8) {$D$};
				\draw[dashed] (0,4)--(4.5,4)  (1.5,0)--(1.5,8);
				\draw[thick] (4.5,-2)--(4.5,11);
			\end{tikzpicture}
		\end{center}
		Miền nghiệm của hệ bất phương trình $\heva{&12x+9y\leq 900\\&0\leq x\leq 45\\&0\leq y\leq 80}$ là miền ngũ giác $OABCD$  trong đó  $O\left(0;0\right)$, $A\left(45;0\right)$,  $B\left(45;40\right)$,  $C\left(15;80\right)$, $D\left(0;80\right)$. \\     
		Tại $O\left(0;0\right)$ , ta có $P=250\,000\cdot0+180\,000\cdot0=0$.\\  
		Tại $A\left(45;0\right)$ ta có  $P=250\,000\cdot45+180\,000\cdot0=11250000$.\\
		Tại  $B\left(45;40\right)$ ta có $P=250\,000\cdot45+180\,000\cdot40=18\,450\,000$. \\
		Tại  $C\left(15;80\right)$ ta có  $P=250\,000\cdot15+180\,000\cdot80=18\,150\,000$. \\
		Tại $D\left(0;80\right)$  ta có  $P=250\,000\cdot0+180\,000\cdot80=14\,400\,000$.\\ 
		Ta thấy biểu thức $P=250\,000x+180\,000y$   đạt giác trị lớn nhất khi $x_0=45, \, y_0=45$.\\
		Vậy $T=x_0+2y_0=125$.		
	}
\end{ex}

\begin{ex}%[0D3N1-3]
	Gia đình chị Minh dự định trồng rau và hoa trên một mảnh đất có diện tích  $8$ ha. Nếu trồng $1$  ha rau thì cần $20$  ngày công và thu lợi $3$ triệu đồng. Nếu trồng $1$  ha hoa thì cần $30$  ngày công và thu lợi $4$  triệu đồng. Biết rằng, gia đình chị Minh chỉ có thể sử dụng không quá $180$  ngày công cho công việc trồng rau và hoa. Hỏi từ việc trồng rau và hoa nói trên, chị Minh có thể thu về lợi nhuận cao nhất là bao nhiêu triệu đồng.
	\shortans{$26$}
	\loigiai{
		Gọi $x, \, y\, \left(x\geq0, \,y\geq0\right)$    lần lượt là số ha đất trồng rau và hoa.\\
		Diện tích đất trồng canh tác không vượt quá $8$  ha nên ta có $x+y\leq8$.\\
		Số ngày công sử dụng không vượt quá $180$ ngày nên $2x+3y\leq180$.\\
		Từ đó, ta có hệ bất phương trình $\heva{&x+y\leq8\\&2x+3y\leq180\\&x\geq 0\\&y\geq 0.}$\\ 
		Ta cần tìm  $x, \, y$   sao cho $T\left(x,y\right)=3x+4y$  lớn nhất.\\
		Miền nghiệm của hệ  được biểu diễn như sau:
		\begin{center}
			\begin{tikzpicture}[scale=0.6,thick,>=stealth']
				\draw[->] (-3,0) -- (12,0)node[above]{$x$};
				\foreach \x in {1,2,3,4,5,6,7,8,9,10,11}
				\draw[shift={(\x,0)},color=black] (0pt,2pt) -- (0pt,-2pt) node[below] {\footnotesize $\x$};
				\draw[->,color=black] (0,-2) -- (0,9)node[right]{$y$};
				\foreach \y in {1,2,3,4,5,6,7,8,9}
				\draw[shift={(0,\y)},color=black] (2pt,0pt) -- (-2pt,0pt) node[left] {\footnotesize $\y$};
				\node[below left] at (0,0) {$O$};
				\clip(-3,-2) rectangle (12,9);
				%\fill[pattern=north east lines] (-0.15,14.3) -- (-1,14.3) -- (-1,-1) -- (7.5,-1)-- cycle;
				\draw[line width=1.2pt,smooth,samples=100,domain=-3:9] plot(\x,{-(\x)+8});
				%\fill[pattern=north east lines] (-1,6.4)--(-1,-1)--(15.3,-1)--(15.3,-0.12)-- cycle;
				\draw[line width=1.2pt,smooth,samples=100,domain=-1:12] plot(\x,{-(2/3)*(\x)+6});			
				\fill[pattern=north east lines] (0,0)--(0,6) --(6,2)--(8,0)-- cycle;
				\node[above] at (8,0) {$A$};
				\node[right] at (6,2) {$B$};
				\node[right] at (0,6) {$C$};			
			\end{tikzpicture}
		\end{center}
		Miền nghiệm của hệ bất phương trình trên là miền trong của tứ giác $OABC$, kể cả các  cạnh của tứ giác đó, với  $O\left(0;0\right)$, $A\left(8;0\right)$,  $B\left(6;2\right)$,  $C\left(0;6\right)$. \\
		Tại $O\left(0;0\right)$, ta có $T=3\cdot0+4\cdot0=0$.\\ 
		Tại $A\left(8;0\right)$, ta có  $T=3\cdot8+4\cdot0=24$.\\ 
		Tại $B\left(6;2\right)$, ta có  $T=3\cdot6+4\cdot2=26$.\\
		Tại $C\left(0;6\right)$, ta có $T=3\cdot0+4\cdot6=24$.\\ 
		Vậy số lợi nhuận cao nhất mà gia đình chị Minh thu được từ trồng rau và hoa là $26$ triệu đồng.	
	}
\end{ex}
 \Closesolutionfile{ans}
\indapan{6}{ans/ans-0-De2-KQ}		